% =====================================================================
%  TFM — Máster Universitario en Inteligencia Artificial
%  Autor: Daniel Gómez Arranz  ·  Convocatoria: Septiembre 2026
%  Compilar:  pdflatex main → biber main → pdflatex main → pdflatex main
%  (o subir la carpeta a Overleaf y compilar con biber).
% =====================================================================
\documentclass[12pt,a4paper]{report}
\usepackage[utf8]{inputenc}
\usepackage[T1]{fontenc}
\usepackage[spanish,es-tabla]{babel}
\usepackage{graphicx,booktabs,amsmath,amssymb,siunitx,array}
\usepackage[a4paper,top=2.5cm,bottom=2.5cm,left=3cm,right=3cm]{geometry}
\usepackage{setspace}
\usepackage{caption}
\usepackage{csquotes}
\usepackage{tikz}
\usetikzlibrary{arrows.meta,positioning,fit,backgrounds}
\usepackage{placeins}
\usepackage{hyperref}
\hypersetup{colorlinks=true,linkcolor=black,citecolor=blue,urlcolor=blue}

% --- Color de títulos según la plantilla oficial UAX (azul #2F5496) ---
\usepackage{xcolor}
\usepackage{titlesec}
\definecolor{uaxblue}{HTML}{2F5496}
\titleformat{\chapter}[display]
  {\normalfont\Huge\bfseries\color{uaxblue}}{\chaptertitlename\ \thechapter}{20pt}
  {\Huge\raggedright\hyphenpenalty=10000\exhyphenpenalty=10000}
\titleformat{\section}
  {\normalfont\Large\bfseries\color{uaxblue}}{\thesection}{1em}{}
\titleformat{\subsection}
  {\normalfont\large\bfseries\color{uaxblue}}{\thesubsection}{1em}{}

% --- Bibliografía APA 6ª (requisito UAX) ---
\usepackage[style=apa,backend=biber]{biblatex}
\DeclareLanguageMapping{spanish}{spanish-apa}
\addbibresource{refs.bib}

\graphicspath{{figuras/}}
\onehalfspacing
% El estilo "Texto TFM" de la plantilla Word (alias del estilo Normal) no usa sangría de
% primera línea y añade 24pt de espacio antes y después de cada párrafo; replicado aquí.
\setlength{\parindent}{0pt}
\setlength{\parskip}{10pt}

% --- Encabezado y pie de página según la plantilla oficial UAX ---
\usepackage{fancyhdr}
\newcommand{\tfmautor}{Gómez Arranz, Daniel}
\newcommand{\tfmtitulocorto}{Fiabilidad de informes radiológicos generados por IA}
\newcommand{\tfmheaderfooter}{%
  \fancyhf{}
  \renewcommand{\headrulewidth}{0.4pt}
  \renewcommand{\footrulewidth}{0.4pt}
  \fancyhead[L]{\small\tfmautor}
  \fancyhead[R]{\raisebox{-0.28\height}{\includegraphics[width=2.6cm]{uax_business_tech_wordmark.png}}}
  \fancyfoot[L]{\small\tfmtitulocorto}
  \fancyfoot[R]{\small\thepage}
}
\pagestyle{fancy}
\tfmheaderfooter
\fancypagestyle{plain}{\tfmheaderfooter}
\setlength{\headheight}{32pt}
\setlength{\headsep}{28pt}
\setlength{\footskip}{40pt}

\begin{document}

% Portada UAX — MUIA (replica la de la plantilla oficial "25-26 Plantilla TFM MUIA").
% WIP: confirmar que "Septiembre 2026" (línea 21) sigue siendo la convocatoria correcta.
\begin{titlepage}
\centering
\includegraphics[width=8cm]{uax_business_tech_wordmark.png}\par
\vspace{0.8cm}
{\large UNIVERSIDAD ALFONSO X EL SABIO\\[2pt]
Business Tech\par}
\vspace{0.5cm}
{\large Máster Universitario en Inteligencia Artificial\par}
\vspace{2cm}
{\Large\bfseries TRABAJO DE FIN DE MÁSTER\par}
\vspace{1.5cm}
{\LARGE Fiabilidad de los Informes Radiológicos Generados\\[4pt]
por Modelos Multimodales: Explicabilidad Visual\\[4pt]
frente a Incertidumbre de la Decisión\par}
\vfill
{\large Autor: Daniel Gómez Arranz\par}
\vspace{0.3cm}
{\large Tutor: José Antonio Marcos García\par}
\vspace{1cm}
{\large Septiembre 2026\par}
\end{titlepage}


\chapter*{Agradecimientos}\addcontentsline{toc}{chapter}{Agradecimientos}

A mi tutor, José Antonio Marcos García, por su orientación durante el desarrollo de este trabajo.

A la Universidad Alfonso X el Sabio y al profesorado del Máster Universitario en Inteligencia
Artificial, por la formación recibida a lo largo de este curso.

A mi familia y amigos, por su apoyo durante estos meses de trabajo.

A mi pareja, cuyas conversaciones sobre su experiencia en el grado de Imagen para el Diagnóstico y
Medicina Nuclear fueron el origen de este proyecto: gracias por acompañarme durante todo el proceso
y por permitirme explorar un ámbito de aplicación de la Inteligencia Artificial con el que no estaba
familiarizado.


% RESUMEN (IMRyD, ~300 palabras) + ABSTRACT en inglés.
\chapter*{Resumen}\addcontentsline{toc}{chapter}{Resumen}

La generación automática de informes de radiografía de tórax mediante modelos multimodales de
lenguaje (MLLM) es prometedora, pero su fiabilidad y la fidelidad de sus explicaciones visuales están
poco estudiadas, lo que limita su adopción clínica. Este trabajo evalúa si la explicación visual de
estos modelos (o su propia incertidumbre) permite estimar la fiabilidad de cada informe. Se
implementa un marco de explicabilidad y se valida el anclaje de la atención de forma objetiva contra
las cajas delimitadoras reales de NIH ChestX-ray8. El estudio de fiabilidad se articula como un diseño
\(2\times2\) que cruza dos generadores del estado del arte (LLaVA-Rad y MAIRA-2) con dos dominios
(Open-I, Estados Unidos; y PadChest-GR, España), etiquetando cada informe como correcto, alucinación u
omisión mediante tres fuentes de verdad independientes (diagnóstico estructurado MeSH, hallazgos
presentes/ausentes de radiólogo y juez basado en modelo de lenguaje). Los resultados muestran que la
atención está clínicamente anclada (lift global 3,03), pero la morfología de la explicación no predice
la fiabilidad. Este resultado se confirma incluso con las cajas de anclaje entrenadas del propio
modelo, que localiza sus alucinaciones de forma plausible. En cambio, la incertidumbre de la decisión, medida
como la probabilidad mínima de contenido, detecta las alucinaciones con AUC 0,90 (Open-I) y se
transfiere de forma robusta a otro modelo y a otro dominio (AUC 0,80--0,81), salvo en un generador
débil fuera de distribución. La señal es asimétrica: detecta las fabricaciones, pero no las omisiones,
que son errores confiados de sensibilidad. Una vez calibrada, la señal se traduce en una confianza por
informe con bandas de riesgo (``semáforo'') y en una demostración interactiva que integra informe,
anclaje y confianza. El hallazgo más relevante del trabajo es que, en los modelos y dominios
evaluados, la detección de alucinaciones mediante incertidumbre generativa resulta una señal
de fiabilidad más útil que la morfología de la explicación visual: es cuantificable, se transfiere
entre modelos y dominios, y es apta para el cribado clínico.

\vspace{0.4cm}\noindent\textbf{Palabras clave:} informes radiológicos, modelos multimodales,
explicabilidad, detección de alucinaciones, incertidumbre.

\chapter*{Abstract}\addcontentsline{toc}{chapter}{Abstract}

Automatic chest X-ray report generation with multimodal large language models (MLLM) is promising, but
the reliability of these models and the faithfulness of their visual explanations remain understudied,
hindering clinical adoption. This work assesses whether a model's visual explanation (or its own
uncertainty) can estimate the reliability of each report. An explainability framework is implemented
and the grounding of the model's attention is objectively validated against the real bounding boxes of
NIH ChestX-ray8. The reliability study follows a \(2\times2\) design crossing two state-of-the-art
generators (LLaVA-Rad and MAIRA-2) with two domains (Open-I, United States; and PadChest-GR, Spain),
labelling each report as correct, hallucination or omission through three independent ground-truth
sources (structured MeSH diagnosis, radiologist present/absent findings, and a language-model judge).
Results show that attention is clinically grounded (global lift 3.03), yet the morphology of the
explanation does not predict reliability. This finding is confirmed even with the model's own trained
grounding boxes, which plausibly localise its hallucinations. Instead, decision uncertainty, measured
as the minimum content-token probability, detects hallucinations with AUC 0.90 (Open-I) and transfers
robustly across model and domain (AUC 0.80--0.81), except for a weak generator out of distribution.
The signal is asymmetric: it detects fabrications but not omissions, which are confident sensitivity
errors. Once calibrated, the signal yields a per-report confidence with risk bands (a ``traffic
light'') and an interactive demo integrating report, grounding and confidence. The most relevant
finding of this work is that, across the models and domains evaluated, hallucination detection
via generative uncertainty is a more useful reliability signal than the morphology of the visual
explanation: it is quantifiable, transfers across models and domains, and is suitable for clinical
triage.

\vspace{0.4cm}\noindent\textbf{Keywords:} radiology reporting, multimodal models, explainability,
hallucination detection, uncertainty.


\tableofcontents
\listoffigures
\listoftables

\chapter{Introducción al Trabajo Fin de Máster}
% BORRADOR para editar a tu voz (estilo impersonal UAX). Recomendado: 5-10 páginas.

La interpretación de imágenes médicas es una de las tareas de mayor valor y, a la vez, de mayor
carga asistencial dentro del sistema sanitario. La radiografía de tórax (CXR) es la prueba de
imagen más solicitada del mundo, y la redacción de su informe consume una parte considerable del
tiempo del radiólogo. En este contexto, la generación automática de informes mediante inteligencia
artificial se presenta como una vía prometedora de asistencia. Sin embargo, para que un sistema de
este tipo sea aceptable en la práctica clínica no basta con que produzca un texto plausible: debe
ser posible saber \textit{cuándo fiarse} de él y \textit{por qué} concluye lo que concluye. Este
capítulo introduce la motivación del trabajo, el problema que aborda y la organización de la
memoria.

\section{Justificación del proyecto realizado}

Mi pareja se encontraba realizando un grado de Imagen para el Diagnóstico y Medicina Nuclear, y debido
a esto tuve la oportunidad de hablar con algunos de los profesores del centro sobre cómo la tecnología
y, en concreto, la inteligencia artificial estaban teniendo un impacto importante en la profesión. De
ahí, la conversación derivó a los principales problemas que tenían con ella (sobre todo la generativa,
no la clasificadora clásica, que lleva muchísimo tiempo en la industria y va por una rama completamente
distinta a la que usa VLMs), y el tema que más se repetía era la falta de explicabilidad, lo ``caja
negra'' que es, que impide a los radiólogos comprender el porqué del diagnóstico, ya que, a diferencia
de los modelos clasificadores, que simplemente clasifican una patología (de manera correcta o
incorrecta), los VLMs dan una explicación de lo que ven en la imagen para apoyar el diagnóstico, pero,
muchas veces, esto resulta en alucinaciones, que producen el efecto contrario al que se quiere generar:
en este caso, agilizar y facilitar el diagnóstico mediante estas radiografías, que, al tener estas
``alucinaciones'', o errores de interpretación, hacen que se dedique más tiempo a la revisión. Esto hizo
que mi planteamiento inicial fuese intentar encontrar una manera de ``explicar'' el diagnóstico del
modelo mediante mapas de calor que interpretasen la ``atención'' del modelo hacia distintas zonas de la
imagen. Si bien, por las propias circunstancias del proyecto, se acabó derivando gran parte del esfuerzo
a otras áreas, la intención seguía siendo la misma: facilitar la comprensión y aumentar la confianza de
los profesionales sanitarios en los VLMs mediante explicabilidad.

Desde el punto de vista de su utilidad, el proyecto se sitúa en la intersección de tres
necesidades actuales. En el ámbito \textbf{clínico}, la presión asistencial y la escasez de
radiólogos hacen atractiva cualquier herramienta que agilice el flujo de trabajo sin comprometer
la seguridad del paciente. En el ámbito \textbf{científico-técnico}, los modelos multimodales de
lenguaje (MLLM) han alcanzado una notable competencia en tareas de imagen y texto, pero su
aplicación a dominios de alta responsabilidad exige garantías de fiabilidad y trazabilidad de las
que todavía carecen. Y en el ámbito \textbf{ético y regulatorio}, la interpretabilidad de los
sistemas de IA en salud no es una característica opcional sino un requisito: un modelo cuyas
decisiones no pueden auditarse difícilmente será adoptado. El presente trabajo se orienta a este
espacio, estudiando no solo la generación de informes sino, sobre todo, la \textbf{fidelidad y el
anclaje clínico de sus explicaciones} y la posibilidad de estimar la \textbf{confianza} en cada
informe generado.

\section{Presentación del proyecto y sus contenidos}
El problema que se aborda puede enunciarse de forma simple: un sistema que genera informes
radiológicos produce, para cada afirmación, una explicación visual (un mapa de calor que señala la
región de la imagen que la motivó); la cuestión es si esa explicación es \textit{fiel} al
razonamiento del modelo, si está \textit{anclada} a la anatomía correcta y, en última instancia,
si permite \textbf{detectar cuándo el modelo se equivoca}.

Considérese, a modo de ilustración, un sistema que informa ``cardiomegalia'' sobre una radiografía
de tórax normal. Si junto al texto se muestra un mapa de calor razonablemente centrado sobre la
silueta cardíaca, el radiólogo dispone de una pista de \emph{dónde} ha mirado el modelo, pero no de
ninguna garantía de que ese mismo patrón de atención no acompañe también a sus aciertos: una
explicación visualmente plausible puede convivir con un contenido incorrecto. La pregunta que
articula el trabajo es, precisamente, si existe alguna señal —en la propia explicación o en la
confianza interna del modelo— capaz de distinguir este caso del de un informe verdaderamente
correcto, antes de que el error llegue al informe final.

Para responderla, el trabajo contribuye en cuatro frentes. En primer lugar, se desarrolla un
\textbf{generador propio} encoder--decoder multimodal, en el que la restricción de datos se
convierte en una decisión metodológica mediante la destilación de conocimiento desde modelos
preentrenados libres. En segundo lugar, se construye un \textbf{marco de explicabilidad} que
evalúa dos ejes: la fidelidad de las explicaciones y su anclaje clínico. En tercer lugar, dado que
un modelo de escala reducida limita el alcance del análisis, se \textbf{eleva el estudio a un MLLM
del estado del arte} (LLaVA-Rad) y se valida su anclaje de forma \textbf{objetiva} contra
anotaciones reales de localización. Por último, se investiga si la explicación —o la propia
incertidumbre del modelo— permite \textbf{detectar alucinaciones}, y se comprueba que la señal de
confianza resultante es \textbf{agnóstica al modelo y robusta al dominio} mediante una réplica en un
segundo generador (MAIRA-2) y un segundo dominio (PadChest-GR); la señal se \textbf{calibra} como una
confianza por informe y se materializa en una demostración interactiva que integra informe, anclaje y
fiabilidad.

El interés del problema se ha intensificado con la irrupción de los \textbf{modelos multimodales de
lenguaje} (MLLM), que acoplan un potente codificador visual a un modelo de lenguaje de gran tamaño y
han elevado de forma notable la calidad del texto generado a partir de imágenes. Su aplicación a la
radiografía de tórax promete asistir en la redacción del informe, agilizando el flujo de trabajo. Sin
embargo, esta misma potencia trae consigo un riesgo característico: la \textbf{alucinación}, es decir,
la generación de afirmaciones plausibles pero no sustentadas por la imagen. En un dominio de alta
responsabilidad como el diagnóstico por imagen, una alucinación no detectada puede tener consecuencias
graves, y una omisión —no reportar un hallazgo presente— otras distintas pero igualmente serias. Por
ello, el problema no es solo generar buen texto, sino \textbf{saber cuándo el texto generado es
fiable}, y hacerlo de una forma que el clínico pueda interpretar y en la que pueda confiar. Es en esta
intersección —entre la capacidad generativa de los MLLM y la exigencia de fiabilidad y transparencia
del entorno clínico— donde se inscribe el presente trabajo.

\section{Preguntas de investigación}
El problema anterior se concreta en cuatro preguntas de investigación que vertebran el trabajo:
\begin{itemize}
  \item \textbf{PI1.} ¿Está la atención de un MLLM de radiografía de tórax \emph{anclada} sobre la
  región anatómica de la lesión, contrastada con anotaciones objetivas de localización?
  \item \textbf{PI2.} ¿Permite la \emph{morfología} de la explicación (la forma del mapa de calor, o
  las cajas de anclaje del propio modelo) detectar cuándo el modelo se equivoca?
  \item \textbf{PI3.} ¿Permite la \emph{incertidumbre} de la decisión del modelo detectar sus errores,
  y con qué asimetría entre alucinaciones y omisiones?
  \item \textbf{PI4.} ¿Es esa señal de fiabilidad \emph{agnóstica al modelo y robusta al dominio}, es
  decir, se transfiere a otro generador y a otra población?
\end{itemize}

\clearpage
\section{Contribuciones}
Las principales contribuciones del trabajo son las siguientes. Primera, una \textbf{auditoría
objetiva del anclaje} de la atención de un MLLM de CXR frente a cajas delimitadoras reales. Segunda,
la constatación —por triple vía, incluida la de las cajas entrenadas del propio modelo— de que la
\textbf{morfología de la explicación no predice la fiabilidad}, así como la identificación de varios
\emph{artefactos} metodológicos que solo un control cuidadoso de confusores reveló. Tercera, la
\textbf{disociación} entre errores detectables por incertidumbre (alucinaciones) y errores confiados
(omisiones), y la propuesta de una señal de confianza \textbf{simple, agnóstica y disponible en
inferencia}. Cuarta, una \textbf{validación cruzada} de dicha señal en un diseño de dos generadores
por dos dominios, que demuestra su robustez. Y quinta, su \textbf{calibración} y su traducción en una
demostración interactiva que integra informe, anclaje y confianza para el clínico. Como aportación
metodológica menor, se documenta que \emph{excluir el token de apertura del informe} mejora de forma
sistemática la señal de incertidumbre.

La memoria se organiza como sigue. El Capítulo~2 fija el objetivo general y los específicos, y
explica cómo evolucionó el alcance del proyecto hasta llegar a ellos. El Capítulo~3 revisa el estado
del arte en cuatro frentes: generación automática de informes, funcionamiento interno de los modelos
multimodales, explicabilidad de \textit{transformers} y detección de incertidumbre y alucinaciones. El
Capítulo~4 describe en detalle la metodología: la población y la muestra, las hipótesis, las variables
y los instrumentos de medida, la arquitectura del generador propio y el diseño experimental. El
Capítulo~5 presenta los resultados siguiendo el orden en que se han enunciado las preguntas de
investigación, y los discute a la luz de la literatura y de sus implicaciones prácticas. El
Capítulo~6, por último, recoge las conclusiones, valora las limitaciones —incluidas las derivadas de
las condiciones materiales del proyecto— y traza las líneas futuras, tanto de investigación como de
mejora del sistema.

\chapter{Objetivos del TFM}
% BORRADOR para editar a tu voz. Recomendado: 1-3 páginas.

En este capítulo se establece el objetivo general del trabajo y los objetivos específicos en que
se desglosa. Los objetivos se redactan en infinitivo y se corresponden con metas concretas
alcanzadas al término del proyecto: uno asociado al marco teórico, otro al desarrollo
metodológico y los restantes a la validación empírica.

\section{Evolución del alcance del proyecto}
El planteamiento inicial situaba la explicabilidad visual —la morfología del mapa de calor de
atención— como la señal candidata para estimar la fiabilidad de un informe. La evidencia recabada
en las primeras fases del trabajo (Capítulo~5, PI2) mostró que dicha morfología no discrimina entre
aciertos y errores, lo que exigió una reformulación del alcance: el foco se desplazó a la
incertidumbre de la decisión como señal alternativa, y el diseño experimental se amplió a una
validación cruzada \(2\times2\) (dos generadores por dos dominios) para descartar que el resultado
negativo de la explicabilidad fuese un artefacto de un único modelo o población. El propio orden de trabajo responde a una decisión deliberada: se optó por construir primero un
\textbf{generador propio}, más pequeño y con destilación de conocimiento, precisamente para entender
\emph{desde dentro} el funcionamiento de estos sistemas y las limitaciones específicas de un modelo
de menor capacidad, antes de dar el salto a un modelo del estado del arte ya entrenado y tratado como
caja cerrada. Esa fase reveló, entre otras lecciones, que las métricas de solapamiento léxico como
ROUGE son inadecuadas para valorar el equilibrio entre acierto y alucinación (Sección~5.2), lo que
motivó abandonarlas en favor de la corrección clínica. Junto con la combinación de cómputo en
hardware de consumo y el acceso restringido a los grandes corpus clínicos (Sección~4.1), estas
lecciones llevaron a elevar el análisis de explicabilidad a un modelo del estado del arte —LLaVA-Rad,
y después MAIRA-2—, ejecutados localmente mediante cuantización en lugar de entrenados desde cero.
Los objetivos que siguen recogen ya el alcance resultante de esta reformulación, no el planteamiento
original.

\section{Objetivo general}
Evaluar la fidelidad y el anclaje clínico de las explicaciones visuales en la generación
automática de informes de radiografía de tórax mediante arquitecturas encoder--decoder
multimodales, y determinar en qué medida dichas explicaciones —o la incertidumbre del propio
modelo— permiten estimar la fiabilidad de cada informe.

\section{Objetivos específicos}
\begin{enumerate}
  \item \textbf{Analizar el estado del arte} en generación automática de informes radiológicos,
  en explicabilidad de modelos transformer y en detección de incertidumbre y alucinaciones,
  identificando el hueco que este trabajo aborda.
  \item \textbf{Diseñar e implementar} un sistema de generación de informes encoder--decoder
  multimodal, incorporando destilación de conocimiento desde modelos preentrenados libres como
  respuesta metodológica a la restricción de datos, junto con un marco reutilizable de
  explicabilidad.
  \item \textbf{Cuantificar de forma objetiva el anclaje clínico} de la atención de un modelo
  multimodal del estado del arte, contrastándola con anotaciones reales de localización
  (bounding boxes de NIH ChestX-ray8).
  \item \textbf{Determinar si la explicabilidad y la incertidumbre} del modelo permiten detectar
  sus errores (alucinaciones y omisiones), sentando las bases de una estimación de confianza por
  informe, agnóstica al modelo empleado.
\end{enumerate}

\chapter{Marco Teórico / Estado del Arte}

Este capítulo revisa el conocimiento científico relevante para el trabajo, organizado en cuatro
áreas: la generación automática de informes radiológicos, la explicabilidad de los modelos
\textit{transformer}, la estimación de incertidumbre y la detección de alucinaciones en modelos
generativos, y los conjuntos de datos y modelos preentrenados del dominio radiológico. El capítulo
concluye con una síntesis que sitúa la aportación del trabajo respecto al estado del arte y delimita
el hueco de investigación que se aborda.

\section{Generación automática de informes radiológicos}

\subsection{De la descripción de imágenes al informe clínico}
La generación de informes a partir de imágenes médicas hereda su formulación del \textit{image
captioning}: producir una secuencia de texto condicionada a una imagen. Las primeras aproximaciones
adoptaron arquitecturas \textbf{codificador--decodificador} en las que una red convolucional
resumía la imagen en un vector de características y una red recurrente (LSTM o GRU) generaba el texto
palabra a palabra. La incorporación de mecanismos de \textbf{atención visual} permitió que cada
palabra generada se condicionara a distintas regiones de la imagen, mejorando tanto la calidad como
la interpretabilidad aparente de la salida.

El informe radiológico presenta, sin embargo, particularidades que lo distinguen de una simple
descripción. Es un texto \textbf{largo, estructurado y altamente estandarizado}, organizado en
secciones (\textit{Findings}, \textit{Impression}) y dominado por un lenguaje clínico repetitivo en
el que la mayoría de estudios son normales. Estas características motivaron arquitecturas específicas:
modelos con \textbf{memoria relacional} y módulos de generación guiada por conocimiento, que buscaban
capturar la estructura del informe y las correlaciones entre hallazgos, y que constituyeron el
estándar antes de la irrupción de los grandes modelos.

\subsection{La era de los modelos multimodales de lenguaje}
La generalización de los \textit{transformers} y el preentrenamiento a gran escala dieron lugar a los
\textbf{modelos multimodales de lenguaje} (MLLM), que acoplan un codificador visual a un modelo de
lenguaje de gran tamaño mediante una capa de proyección, siguiendo el paradigma de LLaVA
\parencite{liu2023llava}. Su adaptación al dominio radiológico ha producido generadores de informes
de calidad notablemente superior. Entre ellos destacan \textbf{LLaVA-Rad} \parencite{llavarad2024},
que combina un codificador biomédico con un decodificador Vicuna afinado sobre datos de tórax; y
\textbf{MAIRA-2} \parencite{maira2}, orientado a la \textit{generación anclada} (\textit{grounded
report generation}): además del texto, produce \textbf{cajas delimitadoras} que localizan cada
hallazgo, entrenado sobre una combinación de corpus (MIMIC-CXR, PadChest y USMix). Otros sistemas
recientes exploran el razonamiento por agentes y la integración de herramientas de detección
especializadas. La Tabla~\ref{tab:modelos} resume las características de los modelos empleados o
relevantes para este trabajo.

\begin{table}[t]\centering
\caption{Modelos multimodales de CXR relevantes. RRG = \textit{radiology report generation}.}
\label{tab:modelos}
\begin{tabular}{lllc}
\toprule
Modelo & Base de lenguaje & Entrenamiento & Anclaje explícito \\ \midrule
LLaVA-Rad & Vicuna-7B & MIMIC-CXR & No \\
MAIRA-2   & Vicuna-7B & MIMIC-CXR, PadChest, USMix & Sí (cajas) \\
CXR-LLAVA & Llama-2-7B & Multi-corpus & No \\ \bottomrule
\end{tabular}
\end{table}

\subsection{El problema de la evaluación}
La evaluación de estos sistemas se apoya tradicionalmente en \textbf{métricas de solapamiento léxico}
heredadas del resumen automático, como ROUGE \parencite{lin2004rouge}, BLEU o METEOR, que comparan el
texto generado con el informe de referencia contando $n$-gramas comunes. Estas métricas presentan una
debilidad grave en radiología: al estar el dominio dominado por informes normales y por un lenguaje
de plantilla, \textbf{un modelo que memoriza la plantilla mayoritaria obtiene puntuaciones altas sin
leer realmente la imagen}. Este trabajo constata empíricamente el fenómeno: una variante mejor
anclada a la imagen puede obtener un ROUGE inferior al de un modelo ciego que reproduce la plantilla.

Por ello, la comunidad ha desarrollado \textbf{métricas de eficacia clínica} que evalúan la
corrección de los hallazgos y no la coincidencia superficial de palabras: clasificadores de
etiquetas radiológicas que extraen los hallazgos de ambos informes y los comparan (obteniendo
precisión, sensibilidad y $F_1$ por patología), y métricas basadas en grafos de conocimiento que
contrastan las entidades clínicas y sus relaciones. Entre las primeras destacan los
\textbf{clasificadores de etiquetas radiológicas} como \textbf{CheXbert} \parencite{smit2020chexbert},
que combina reglas expertas y un modelo BERT afinado para extraer de cada informe un vector de
hallazgos —presente, ausente o incierto para cada patología— y calcular métricas de precisión,
sensibilidad y $F_1$ sobre esa comparación; entre las segundas, las basadas en \textbf{grafos de
conocimiento} como \textbf{RadGraph} \parencite{jain2021radgraph}, que representa las entidades
clínicas de un informe (hallazgos, anatomía, atributos) y sus relaciones en un grafo estructurado, y
cuya comparación entre grafo generado y grafo de referencia (\textit{RadGraph~$F_1$}) penaliza tanto
los hallazgos inventados como los omitidos con mayor granularidad que un simple recuento de
etiquetas. Estas métricas se aproximan mucho más al criterio del radiólogo que el solapamiento léxico,
aunque a costa de depender de un extractor de hallazgos que introduce su propio error —una limitación
que este trabajo evita adoptando, en su lugar, fuentes de verdad ya estructuradas (MeSH, hallazgos de
radiólogo) y un juez de lenguaje sobre la narrativa completa (Capítulo~4)—. La distinción
entre \emph{corrección clínica} y
\emph{similitud textual} es central para este trabajo, pues motiva que la fiabilidad se defina en
términos de acierto sobre hallazgos significativos y no de solapamiento léxico. En consecuencia, la
frontera de la investigación se ha desplazado de \emph{generar mejor} a \emph{saber cuándo fiarse} de
lo generado: la fiabilidad y la explicabilidad.

\section{Anatomía y funcionamiento de un modelo multimodal de lenguaje}
Comprender cómo funciona internamente un MLLM es imprescindible para justificar por qué las señales
empleadas en este trabajo —el mapa de calor de atención y las probabilidades de los tokens— son
precisamente las que son. Esta sección describe sus componentes.

\subsection{El \textit{transformer} y el mecanismo de atención}
La arquitectura \textit{transformer} \parencite{vaswani2017} sustituyó la recurrencia por el
mecanismo de \textbf{auto-atención}: cada elemento de una secuencia calcula una combinación ponderada
del resto, donde los pesos —la \emph{atención}— se obtienen de la similitud entre una \textit{query} y
unas \textit{keys}. Formalmente, $\mathrm{Attention}(Q,K,V)=\mathrm{softmax}(QK^{\top}/\sqrt{d})V$.
Este mecanismo, además de su eficacia, expone una magnitud interpretable —cuánto atiende cada token a
cada otro— que está en el origen del debate sobre si la atención constituye una explicación, y que
este trabajo aprovecha y, al mismo tiempo, somete a validación.

\subsection{El codificador visual}
Para procesar la imagen, los MLLM emplean un \textbf{Vision Transformer} (ViT)
\parencite{dosovitskiy2021}, que divide la imagen en una rejilla de \textbf{parches} (por ejemplo,
$16\times16$ o $14\times14$ píxeles), los proyecta linealmente en \textit{tokens} y les aplica
auto-atención como si fueran palabras. El resultado es una secuencia de \textbf{características por
parche} más un \textit{token} global (CLS). Esta rejilla de parches es la que confiere \emph{estructura
espacial} a la atención: al promediar la atención que los tokens de texto dirigen a los tokens
visuales y reordenarla en la rejilla, se obtiene el mapa de calor sobre la imagen. En los codificadores
empleados en este trabajo, dicha rejilla es de $37\times37$ parches. Codificadores auto-supervisados de
dominio, como RAD-DINO \parencite{perezgarcia2024}, aprenden estas representaciones sin supervisión
textual y mejoran la calidad del anclaje.

\subsection{El conector y el preentrenamiento contrastivo}
El puente entre la imagen y el modelo de lenguaje —el \textbf{conector}— adopta distintas formas: una
simple \textbf{proyección lineal} de las características visuales al espacio del lenguaje (como en
LLaVA \parencite{liu2023llava}), un módulo de consulta aprendido (\textit{Q-Former}) que resume la
imagen en unos pocos tokens (BLIP-2 \parencite{li2023blip2}), o \textbf{atención cruzada con
compuerta} intercalada en el decodificador (Flamingo \parencite{alayrac2022}). El generador propio de
este trabajo adopta esta última vía. Con frecuencia, el codificador visual se ha alineado previamente
con el texto mediante un \textbf{preentrenamiento contrastivo} de tipo CLIP \parencite{radford2021},
que aproxima en un espacio común los embeddings de imágenes y textos emparejados; la Etapa~1 del
modelo propio reproduce esta idea.

\subsection{El decodificador y la generación}
Una vez inyectadas las características visuales, un \textbf{modelo de lenguaje} decodificador genera el
informe de forma \textbf{autorregresiva}: en cada paso produce una distribución de probabilidad sobre
el vocabulario y emite un token. Este proceso es la segunda fuente de señal que explota el trabajo: la
\textbf{probabilidad} que el modelo asigna a cada token generado —y su entropía— son un subproducto
directo del paso hacia adelante, disponible \emph{sin coste adicional} en cualquier generador. El
entrenamiento suele constar de una fase de alineamiento y una de \textit{instruction tuning} visual
\parencite{liu2023llava}, que enseña al modelo a seguir la instrucción de describir la imagen.

\paragraph{Consecuencia metodológica.} La arquitectura anterior explica por qué las dos señales
centrales del trabajo son \emph{naturales} y \emph{agnósticas al modelo}: la atención de los tokens de
texto hacia los tokens visuales proporciona el \textbf{mapa de calor} (la explicación de \emph{dónde}
mira el modelo), y la distribución de probabilidad de cada token proporciona la \textbf{incertidumbre}
(la señal de \emph{cuánto} fiarse). Ambas se obtienen del mismo paso de inferencia que ya se ejecuta
para generar el informe, sin modelos auxiliares ni cómputo extra.

\section{Por qué abordar el problema desde la fiabilidad y la incertidumbre}
La decisión de centrar el trabajo en la fiabilidad —y no en mejorar la generación— responde a un
diagnóstico del estado del arte. Por un lado, la calidad del texto generado por los MLLM de CXR
\emph{se está saturando}, mientras que el verdadero cuello de botella para su adopción clínica es la
\textbf{confianza}: la tendencia de estos modelos a alucinar sigue siendo la principal barrera
\parencite{ji2023}. Por otro, existen dos grandes familias de señales de fiabilidad: las
\textbf{externas}, que requieren un verificador adicional (un informe de referencia u otro modelo), y
las \textbf{intrínsecas}, que se leen de la propia salida del modelo. Este trabajo prioriza las
segundas por tres razones: son \textbf{agnósticas al modelo} (transfieren a cualquier generador), están
\textbf{disponibles en inferencia} sin coste adicional, y capturan la información de que el propio
modelo dispone sobre su acierto —los modelos de lenguaje, en buena medida, ``saben lo que saben''
\parencite{kadavath2022}—. La cuestión abierta, que el trabajo aborda, es si esa información se lee
mejor en la \emph{forma} de la explicación visual o en la \emph{incertidumbre} de la decisión, y si,
una vez elegida, resulta \textbf{calibrable} \parencite{guo2017} para su uso clínico.

\section{Explicabilidad en modelos \textit{transformer}}

\subsection{¿Es la atención una explicación?}
La explicabilidad de los \textit{transformers} gira en gran medida en torno a si los pesos de
\textbf{atención} constituyen una explicación válida de sus decisiones. El debate quedó planteado por
\textcite{jain2019}, que mostraron que es posible encontrar distribuciones de atención muy distintas
que producen la misma predicción, cuestionando que la atención sea \emph{fiel} al razonamiento del
modelo. La réplica de \textcite{wiegreffe2019} matizó la conclusión: la atención sí porta información
explicativa bajo ciertas condiciones y su utilidad depende de cómo se defina y evalúe una
explicación. Este debate es directamente pertinente para la radiología, donde el mapa de calor de
atención se ofrece con frecuencia como justificación visual del diagnóstico.

\subsection{Métodos de atribución y agregación}
Para agregar la atención a lo largo de las múltiples capas del modelo se propuso el \textit{attention
rollout} \parencite{abnar2020}, que compone las matrices de atención capa a capa teniendo en cuenta
las conexiones residuales. Métodos posteriores como el de \textcite{chefer2021} combinan atención y
gradientes para obtener mapas de relevancia más fieles, propagando la relevancia a través de la red
de forma consistente con las reglas de conservación. En paralelo, la familia de métodos basados en
\textbf{gradientes} produce mapas de saliencia a partir de los gradientes de la salida respecto a las
activaciones intermedias: el \textit{class activation mapping} y su generalización \textbf{Grad-CAM}
ponderan los mapas de características por la importancia del gradiente para obtener una localización
gruesa de la evidencia, mientras que los métodos de \textbf{propagación de relevancia} redistribuyen la
puntuación de salida hacia atrás, capa a capa, hasta la entrada. Cada familia tiene ventajas e
inconvenientes: los mapas de atención son baratos pero de fidelidad discutida; los basados en
gradientes son más costosos y sensibles al ruido; y los de propagación exigen reglas específicas por
tipo de capa. Frente a ellos, \textit{Integrated
Gradients} \parencite{sundararajan2017} mide la contribución causal de cada entrada integrando los
gradientes a lo largo de un camino desde una referencia neutra hasta la entrada real, ofreciendo
\textbf{garantías axiomáticas} (sensibilidad e invariancia de implementación) de las que la atención
carece.

\subsection{Evaluación objetiva de la fidelidad}
La proliferación de métodos hace imprescindible \textbf{evaluar objetivamente} la fidelidad de una
explicación, en lugar de limitarse a comparar mapas entre sí. Las \textbf{métricas de perturbación}
—eliminación e inserción de las regiones más relevantes— y las nociones de
\textit{comprehensiveness} (cuánto cae la confianza al eliminar lo señalado como importante) y
\textit{sufficiency} (cuánta confianza se conserva reteniendo solo lo importante), formalizadas en el
\textit{benchmark} ERASER \parencite{deyoung2020}, permiten cuantificar si la explicación captura de
verdad lo que impulsa la predicción. Este trabajo adopta esta filosofía: la explicación se valida
frente a \textbf{referencias objetivas} —cajas delimitadoras reales y perturbaciones controladas— en
lugar de asumir su fidelidad por inspección visual. Como advertencia transversal, \textcite{ghassemi2021}
argumentan que las explicaciones locales por saliencia en salud pueden generar una
\textbf{falsa sensación de confianza} y fomentar el sesgo de automatización si se presentan como
garantía de fiabilidad sin validación, lo que refuerza la necesidad de distinguir entre \emph{dónde}
mira el modelo y \emph{cuánto} debe fiarse el clínico de su veredicto.

\section{Incertidumbre y detección de alucinaciones}

\subsection{Naturaleza de las alucinaciones}
En los modelos generativos de lenguaje se denomina \textbf{alucinación} a la producción de
afirmaciones no sustentadas por la evidencia. Los estudios de síntesis del fenómeno
\parencite{ji2023} distinguen entre alucinaciones \emph{intrínsecas} (que contradicen la entrada) y
\emph{extrínsecas} (no verificables con ella), y las atribuyen a factores tanto de los datos —sesgos y
ruido en el corpus de entrenamiento— como de la propia decodificación —la exposición a su propia
salida y la tendencia a completar patrones plausibles—. En
la generación de informes radiológicos, esta taxonomía se concreta en dos modos de error
clínicamente opuestos: la \textbf{alucinación} propiamente dicha —afirmar un hallazgo agudo que no
está presente (falso positivo)— y la \textbf{omisión} o \textit{miss} —no reportar un hallazgo agudo
que sí lo está (falso negativo)—. Ambos comprometen la seguridad del paciente de forma distinta y,
como se mostrará, presentan una \textbf{firma de incertidumbre radicalmente diferente}, lo que
constituye uno de los ejes del presente trabajo.

\subsection{Cuantificación de la incertidumbre}
La detección de alucinaciones se ha vinculado más a la \textbf{incertidumbre del propio modelo} que a
sus explicaciones visuales. Las señales más directas son de \textbf{nivel de token}: la probabilidad
asignada por el modelo a cada token generado y la \textbf{entropía predictiva} de la distribución en
cada paso. La probabilidad mínima a lo largo de la secuencia (el token menos seguro) y la fracción de
tokens de baja probabilidad son resúmenes habituales. Estas señales presentan, no obstante,
\textbf{confusores} conocidos —la longitud del texto, la posición del token o su función sintáctica—
que este trabajo controla explícitamente. Un refinamiento propuesto es la \textbf{entropía semántica}
\parencite{kuhn2023}, que agrupa las generaciones equivalentes en significado antes de medir la
incertidumbre, de modo que la variabilidad puramente léxica no infle la estimación.

\subsection{Auto-consistencia y calibración}
Una vía complementaria y \textbf{agnóstica al modelo} (solo requiere poder muestrear) es la
\textbf{auto-consistencia}: SelfCheckGPT \parencite{manakul2023} detecta alucinaciones sin recursos
externos generando varias muestras y midiendo su acuerdo, bajo la hipótesis de que un contenido
factual tiende a repetirse mientras que uno alucinado varía entre muestras. Esta idea se ha adaptado
específicamente a la generación de informes radiológicos mediante métodos de tipo caja negra que
muestrean varios informes y contrastan la consistencia de cada afirmación, como RadFlag
\parencite{radflag2024}, que combina muestreo y un juez de lenguaje para marcar las afirmaciones poco
respaldadas. El presente trabajo se relaciona con esta línea pero se diferencia de ella en que, además
de la consistencia, estudia la señal de probabilidad de token —más barata, pues no exige múltiples
generaciones—, la caracteriza frente a las omisiones y la valida de forma cruzada entre modelos y
dominios. Finalmente, para que una
señal de confianza sea útil no basta con que discrimine: debe estar \textbf{calibrada}, es decir, que
una confianza estimada del 20\,\% corresponda a un 20\,\% de aciertos observados. La calibración se
evalúa con curvas de fiabilidad y con el error de calibración esperado, y se corrige con métodos como
la regresión de Platt o la regresión isotónica. La literatura advierte, además, de que el ajuste por
refuerzo con retroalimentación humana tiende a producir modelos \textbf{sobre-confiados}, lo que
puede degradar precisamente las señales de probabilidad en modelos muy afinados. Estas ideas motivan
que en este trabajo se evalúen la probabilidad de los tokens y la variabilidad entre muestras como
señales candidatas, que se controle su calibración y que se declaren sus límites. La
Tabla~\ref{tab:senales_lit} resume las familias de señales revisadas y su coste.

\begin{table}[t]\centering
\caption{Familias de señales de incertidumbre revisadas en la literatura. Coste en inferencia: sin
generaciones adicionales (bajo) o con varias muestras extra (alto).}
\label{tab:senales_lit}
\begin{tabular}{p{4.3cm}lp{2.9cm}p{3.3cm}}
\toprule
Señal & Nivel & Coste & Referencia \\ \midrule
Probabilidad de token (mín./media) & Token & Bajo & Este trabajo \\
Entropía predictiva & Token & Bajo & --- \\
Entropía semántica & Secuencia & Alto (agrupa muestras) & \parencite{kuhn2023} \\
Auto-consistencia (SelfCheckGPT) & Secuencia & Alto ($K$ muestras) & \parencite{manakul2023} \\
Auto-consistencia + juez (RadFlag) & Secuencia & Alto ($K$ muestras + juez) & \parencite{radflag2024} \\
\bottomrule
\end{tabular}
\end{table}

\section{Datos y modelos preentrenados en radiología}

\subsection{Conjuntos de datos}
La disponibilidad de datos condiciona fuertemente el dominio. El conjunto \textbf{NIH ChestX-ray8}
\parencite{wang2017} aporta clasificación multi-etiqueta y, de forma parcial, \textbf{localización
mediante cajas delimitadoras} para ocho patologías, lo que lo hace idóneo como referencia objetiva de
anclaje. \textbf{Open-I / Indiana University X-ray} \parencite{demnerfushman2016} es prácticamente el
único conjunto de acceso libre con \textbf{informes narrativos completos} y, además, con diagnóstico
estructurado (términos MeSH), lo que permite un etiquetado objetivo del error. Corpus de mayor tamaño
como \textbf{MIMIC-CXR} ofrecen cientos de miles de informes pero requieren credenciales
restringidas. \textbf{PadChest} y su subconjunto anclado \textbf{PadChest-GR} \parencite{padchestgr2024}
aportan un dominio distinto (población española), anotaciones de hallazgos \textbf{presentes y
ausentes} por radiólogo y cajas, resultando idóneos tanto como segunda fuente de verdad objetiva
—independiente del juez basado en modelo de lenguaje— como para evaluar la generalización fuera de
distribución.

\subsection{Codificadores y destilación de conocimiento}
La restricción de datos narrativos motiva el uso de \textbf{destilación de conocimiento} desde
modelos preentrenados libres. En el eje textual, \textbf{CXR-BERT} \parencite{boecking2022} es un
codificador de lenguaje con semántica clínica, entrenado para alinear texto e imagen radiológicos; en
el eje visual, \textbf{RAD-DINO} \parencite{perezgarcia2024} es un codificador de imagen radiológica
auto-supervisado que aprende representaciones potentes sin depender de supervisión textual, y que
constituye, de hecho, el codificador visual de MAIRA-2. La destilación permite mejorar un modelo
propio de escala reducida sin acceso a datos adicionales, convirtiendo la limitación de datos en una
decisión metodológica deliberada. El panorama de generadores del estado del arte es, además,
crecientemente diverso: junto a los modelos densos de tipo LLaVa han aparecido sistemas de
\textit{grounded reporting} que producen cajas, modelos de propósito médico general y enfoques de
razonamiento por agentes que combinan varios modelos especializados. Esta diversidad refuerza el
interés de disponer de una señal de fiabilidad \textbf{agnóstica al generador}, que pueda aplicarse a
cualquiera de ellos sin depender de su arquitectura interna, como la que este trabajo propone.

\section{Síntesis y hueco de investigación}
De la revisión anterior se desprende un panorama consistente. Los MLLM de CXR del estado del arte
generan texto de calidad y muestran alta especificidad, pero adolecen de dos carencias: una
\textbf{sensibilidad limitada} —tienden a omitir hallazgos, con una especificidad notablemente
superior a su sensibilidad ante la patología \parencite{frontiers2025}— y, sobre todo, una
\textbf{ausencia de explicabilidad validada}, como se ha discutido en las Secciones~3.2 y~3.4. Al
mismo tiempo, la literatura de explicabilidad
advierte de que las explicaciones visuales pueden ser \textbf{engañosas} si se asumen fiables sin
validación \parencite{ghassemi2021}, y la literatura de incertidumbre ofrece señales prometedoras
—probabilidad de token, entropía semántica, auto-consistencia— pero desarrolladas, en su mayoría,
sobre modelos de lenguaje generales y sin una validación cruzada entre modelos y dominios en el
ámbito radiológico. Este trabajo se sitúa precisamente en ese hueco: en lugar de proponer otro
generador, \textbf{audita objetivamente} si la explicación de un MLLM de CXR está anclada, comprueba
si su morfología permite detectar los errores, y determina dónde reside realmente la señal de
fiabilidad, validándola de forma cruzada sobre dos modelos y dos dominios y traduciéndola a una
confianza calibrada por informe.

\chapter{Marco Metodológico}
% BORRADOR para editar a tu voz (estilo impersonal UAX). Recomendado: 5-15 páginas.

En este capítulo se describe el método seguido para desarrollar y evaluar el trabajo. Se
detallan la población y la muestra empleadas, los objetivos de investigación y sus hipótesis,
las variables medidas y los instrumentos de evaluación, el diseño experimental y, finalmente,
el modelo de análisis de datos. El enfoque combina el desarrollo de un sistema de generación de
informes con un marco de evaluación de la explicabilidad, y su validación tanto sobre un modelo
propio como sobre un modelo multimodal del estado del arte.

\section{Población y muestra}
La población objeto de estudio son los estudios de radiografía de tórax (CXR) acompañados de su
informe clínico. Se han utilizado dos conjuntos de datos de acceso libre, seleccionados por su
disponibilidad sin restricciones frente a alternativas como MIMIC-CXR, cuyo acceso exige
credenciales de PhysioNet no disponibles para este trabajo.

El conjunto principal es \textbf{Open-I / Indiana University CXR} \parencite{demnerfushman2016},
que reúne 3.955 estudios y 7.470 imágenes (vistas frontal y lateral) con informes narrativos
estructurados en secciones \textit{FINDINGS} e \textit{IMPRESSION}. La partición en
entrenamiento, validación y prueba (70/15/15) se realiza \textbf{por identificador de estudio},
de modo que ambas vistas de un mismo paciente permanecen en la misma partición y se evita la
fuga de información entre conjuntos.

El segundo conjunto es \textbf{NIH ChestX-ray8} \parencite{wang2017}, del que se emplean las 984
anotaciones con \textbf{caja delimitadora (bounding box)} correspondientes a ocho patologías
(atelectasia, cardiomegalia, derrame, infiltrado, masa, nódulo, neumonía y neumotórax). Estas
cajas constituyen la \textbf{referencia objetiva de localización} que permite evaluar el anclaje
clínico de la atención sin depender de un juicio subjetivo. La restricción de datos narrativos
—Open-I es prácticamente el único conjunto libre con informes completos— motiva que la mejora
del modelo propio se aborde mediante \textbf{destilación de conocimiento} desde modelos
preentrenados libres, convirtiendo la limitación en una decisión metodológica.

El tercer conjunto es \textbf{PadChest-GR} \parencite{padchestgr2024}, subconjunto anclado del corpus
español PadChest, del que se emplea una muestra de estudios con \textbf{hallazgos presentes y ausentes
anotados por radiólogo} y sus cajas. PadChest-GR cumple una doble función: proporciona una
\textbf{segunda fuente de verdad objetiva} —independiente del juez basado en modelo de lenguaje— y
constituye un \textbf{dominio distinto} (población española) respecto a Open-I (Estados Unidos), lo
que permite evaluar la generalización \emph{fuera de distribución}. El cruce de ambos dominios con
dos generadores da lugar al diseño \(2\times2\) que vertebra la validación.

En cuanto al \textbf{detalle y el preprocesado}, Open-I reúne 3.955 estudios y 7.470 imágenes
(proyecciones frontal y lateral); su partición 70/15/15 se realiza por identificador de estudio para
que ambas vistas de un paciente permanezcan juntas y no se produzca fuga de información. De NIH se
emplean las 984 anotaciones con caja de ocho patologías. Del subconjunto de PadChest-GR analizado
—382 estudios, de los cuales 120 son normales y 262 presentan algún hallazgo— se extraen los hallazgos
presentes y ausentes anotados por radiólogo. En todos los casos, las imágenes se convierten a escala de
grises de tres canales, se redimensionan a la resolución de entrada del codificador correspondiente
(224 píxeles para el \textit{Vision Transformer}, 518 para RAD-DINO) y se normalizan con las
estadísticas del preentrenamiento. El texto de referencia se limpia de secciones administrativas,
conservando \textit{Findings} e \textit{Impression}.

\section{Objetivos de investigación}
La investigación empírica se organiza en torno a tres hipótesis:
\begin{itemize}
  \item \textbf{H1.} La atención de un modelo multimodal de CXR está \textit{anclada
  clínicamente}, es decir, se concentra sobre la región anatómica de la lesión.
  \item \textbf{H2.} La \textit{forma} del mapa de calor (su morfología: concentración,
  dispersión, etc.) permite detectar cuándo el modelo se equivoca (alucinación o error).
  \item \textbf{H3.} La \textit{incertidumbre generativa} del modelo (las probabilidades de su
  propia salida) permite detectar dichos errores.
  \item \textbf{H4.} Dicha señal de incertidumbre es \textit{agnóstica al modelo y robusta al
  dominio}: se transfiere a otro generador y a otra población distinta de la de entrenamiento.
\end{itemize}
Como se mostrará, H1 se confirma y H3 se sostiene, mientras que H2 se refuta de forma robusta: la
fiabilidad no reside en la morfología de la explicación sino en la incertidumbre de la decisión del
modelo. H4 se confirma con un matiz relevante: la señal se transfiere entre modelos y se mantiene
fuera de distribución para un generador de calidad suficiente, mientras que un generador débil la
pierde fuera de distribución al recaer en un informe-plantilla emitido con confianza.

\section{Variables intervinientes e instrumentos de evaluación}
Las \textbf{variables} evaluadas se agrupan en tres familias. (i) \textit{Morfología del mapa de
calor}: concentración, dispersión (entropía, índice de Gini), agrupamiento y autocorrelación
espacial. (ii) \textit{Incertidumbre generativa}: probabilidad del token menos seguro
(\texttt{min\_prob}) y entropía predictiva por token. Una decisión de diseño relevante es emplear la
\textbf{probabilidad mínima de \emph{contenido}} —excluyendo el token de apertura del informe, cuya
incertidumbre es de estilo y no de contenido—, que mejora sustancialmente la señal (véase
Cap.~5). (iii) \textit{Dependencia de la imagen y consistencia}: fracción de atención dirigida a los
tokens visuales frente al texto, persistencia del hallazgo ante una imagen en blanco
(\emph{blank\_sim}) y variabilidad entre muestras (auto-consistencia tipo SelfCheckGPT). La
\emph{blank\_sim} resulta además un indicador de la calidad del generador: cuanto más depende de la
imagen (menor similitud con la salida en blanco), mejor anclado está el modelo.

Los \textbf{instrumentos} de evaluación son: las métricas de anclaje \textit{pointing-game},
\textit{energy-in-bbox} y \textit{lift} (energía dentro de la caja dividida por su área, es decir,
la mejora sobre el azar); las métricas de discriminación (área bajo la curva ROC, AUC) con
\textit{test de permutación} e \textbf{intervalos de confianza por \emph{bootstrap}}; y el
etiquetado del error, para el que se emplean \textbf{tres fuentes de verdad independientes} a fin de
evitar la circularidad de un único juez: (a) el \textbf{diagnóstico estructurado oficial} (términos
MeSH de Open-I; hallazgos presentes/ausentes anotados por radiólogo en PadChest-GR), objetivo y
reproducible; y (b) el \textbf{juez basado en un modelo de lenguaje} (metodología tipo G-Eval), que
compara la \textbf{narrativa completa} generada con la referencia y la clasifica como correcta,
alucinada u omitida según una rúbrica predefinida. La concordancia entre las tres fuentes (por
ejemplo, AUC 0,87 con juez, 0,82 con MeSH oficial sobre Open-I) sustenta la fiabilidad del
etiquetado. El juez se ejecuta con escritura incremental y de forma reproducible sobre cada informe.

\begin{table}[t]\centering
\caption{Conjuntos de datos empleados y su función en el trabajo.}
\label{tab:datos}
\begin{tabular}{llll}
\toprule
Conjunto & Contenido & Función & Verdad \\ \midrule
Open-I & Informes + MeSH & Estudio principal (in-dist.) & MeSH + juez \\
NIH ChestX-ray8 & Cajas (8 patologías) & Anclaje objetivo & Cajas \\
PadChest-GR & Presentes/ausentes + cajas & Réplica (dominio, OOD) & Radiólogo + juez \\ \bottomrule
\end{tabular}
\end{table}

\section{El marco de explicabilidad}
El marco de explicabilidad implementado evalúa dos ejes complementarios: \emph{dónde} mira el modelo y
si \emph{de verdad} depende de ello. El \textbf{mapa de calor} se obtiene agregando la atención que
los tokens del informe dirigen a los tokens visuales del codificador. Para los modelos de tipo LLaVa,
los tokens visuales forman una rejilla de $37\times37$ (1369 parches); se promedia la atención sobre
capas y cabezas, se restringe a los tokens visuales y se reordena en la rejilla, obteniendo un mapa
espacial que se superpone a la imagen. Se descarta la contribución de los tokens ``sumidero''
(posiciones de atención residual sin significado espacial). Sobre este mapa se calculan las
\textbf{métricas de anclaje} frente a las cajas reales de NIH: el \textit{energy-in-bbox} (fracción de
energía de atención dentro de la caja), el \textit{lift} (dicha fracción normalizada por el área,
es decir, la mejora sobre el azar) y el \textit{pointing-game} (si el máximo del mapa cae dentro de la
caja).

La \textbf{fidelidad} de la explicación se evalúa por \textbf{perturbación}: se ocluye la región de
máxima atención y se mide la caída de la log-probabilidad del hallazgo (\emph{occ\_drop}). Una
explicación fiel debería producir una caída apreciable —si la región atendida realmente sustenta la
afirmación—; su ausencia indicaría que la atención no es causalmente responsable de la salida. Como se
detalla en el capítulo de resultados, esta y otras métricas de forma \emph{no} logran predecir la
fiabilidad, lo que motiva el giro hacia la incertidumbre de la decisión.

\section{Arquitectura del generador propio}
El generador propio es un modelo \textbf{codificador--atención cruzada--decodificador} para la
generación de informes de tórax. El \textbf{codificador visual} es un \textit{Vision Transformer}
(ViT) preentrenado que produce, para cada imagen, un embedding global (CLS) y una rejilla de
características por parche; en las variantes más avanzadas se sustituye por \textbf{RAD-DINO}
\parencite{perezgarcia2024} congelado (518\,px, rejilla de $37\times37$ parches, 768 dimensiones),
aprovechando un codificador de dominio de mayor calidad. El \textbf{decodificador} es un modelo de
lenguaje GPT-2 al que se \textbf{inyecta atención cruzada} hacia las características visuales en cada
bloque. Un aspecto clave de diseño es que dicha atención cruzada se introduce mediante una
\textbf{compuerta inicializada a cero}: al comienzo del entrenamiento la contribución visual es nula y
el decodificador se comporta exactamente como el GPT-2 preentrenado, lo que \emph{preserva} sus pesos
y evita el colapso de la generación; la compuerta aprende progresivamente cuánta información visual
incorporar.

El entrenamiento se organiza en \textbf{dos etapas}. En la \textbf{Etapa 1 (alineamiento
visual--textual)} se entrena únicamente el codificador y unas cabezas de proyección con una pérdida
contrastiva de tipo InfoNCE (CLIP): los embeddings de imágenes e informes emparejados deben acercarse
frente a los no emparejados. En la \textbf{Etapa 2 (ajuste \mbox{generativo})} se entrena el modelo completo
con una pérdida de modelado de lenguaje (entropía cruzada sobre el siguiente token). Opcionalmente,
mediante ajuste eficiente en parámetros (PEFT), se congela el GPT-2 base y se entrenan solo las capas
de atención cruzada, reduciendo los parámetros entrenables de unos 200 a unos 5 millones y mitigando
la ceguera a la imagen.

Ante la restricción de datos narrativos, el modelo incorpora \textbf{destilación de conocimiento}, que
se estudia mediante una ablación progresiva (M0--M3): \textbf{M0}, línea base sin destilación; en el
eje textual, se destila desde \textbf{CXR-BERT} \parencite{boecking2022} —bien como torre de texto
congelada que aporta un embedding clínico, bien como objetivo de destilación que el modelo
\emph{student} aprende a predecir—; y en el eje visual (\textbf{M3}), se destila el embedding CLS del
ViT hacia el de RAD-DINO. Esta ablación permite aislar la contribución de cada fuente de conocimiento
externo a la calidad y al anclaje del modelo resultante. La Figura~\ref{fig:arch} esquematiza la
arquitectura y sus dos etapas de entrenamiento.

\begin{figure}[t]\centering
\begin{tikzpicture}[
  node distance=7mm and 20mm,
  enc/.style={draw, rounded corners, align=center, minimum height=16mm, minimum width=22mm, font=\footnotesize, fill=blue!6},
  dec/.style={draw, rounded corners, align=center, minimum height=16mm, minimum width=26mm, font=\footnotesize, fill=orange!8},
  io/.style={draw, align=center, minimum height=9mm, font=\footnotesize, fill=gray!8},
  flow/.style={-{Latex[length=2mm]}, thick}]
  \node[io] (img) {Radiografía};
  \node[enc, right=of img] (enc) {Codificador visual\\(ViT / RAD-DINO)};
  \node[dec, right=of enc] (dec) {Decodificador GPT-2\\+ atención cruzada\\(compuerta $\to 0$)};
  \node[io, right=of dec] (rep) {Informe};
  \draw[flow] (img) -- (enc);
  \draw[flow] (enc) -- node[above=1.5mm, font=\scriptsize, align=center]{parches\\$37{\times}37$} (dec);
  \draw[flow] (dec) -- (rep);
  \node[font=\scriptsize, below=3mm of enc, text width=30mm, align=center] {\textbf{Etapa 1}: alineamiento contrastivo (InfoNCE)};
  \node[font=\scriptsize, below=3mm of dec, text width=34mm, align=center] {\textbf{Etapa 2}: ajuste generativo (modelado de lenguaje)};
\end{tikzpicture}
\caption{Arquitectura del generador propio: codificador visual, decodificador con atención cruzada
mediante compuerta inicializada a cero, y las dos etapas de entrenamiento (alineamiento y generación).}
\label{fig:arch}
\end{figure}

\section{Definición formal de las señales}
Para garantizar la reproducibilidad, se definen formalmente las señales empleadas. Sea un informe
generado la secuencia de tokens $(w_1,\dots,w_n)$ con probabilidades $p_t = P(w_t \mid w_{<t},
\text{imagen})$ asignadas por el modelo en cada paso $t$.

\paragraph{Incertidumbre generativa.} La señal principal es la \textbf{probabilidad mínima de
contenido}, que excluye el token de apertura (cuya incertidumbre es de estilo y no de contenido):
\begin{equation}
u_{\min} \;=\; -\!\!\min_{2 \le t \le n} p_t .
\end{equation}
Se calcula también la \textbf{entropía predictiva} por token, $H_t = -\sum_{v} p_t(v)\log p_t(v)$, de
la que se toma el máximo $H_{\max}=\max_t H_t$. Para eliminar su confusión con la longitud, la
entropía se \textbf{residualiza} respecto al número de tokens $n$ mediante una regresión lineal
$H_{\max} = \alpha + \beta n + \varepsilon$, tomando el residuo $\varepsilon$ como señal limpia.

\paragraph{Auto-consistencia.} Dadas $K$ muestras $\{s_1,\dots,s_K\}$ generadas por muestreo
estocástico, se mide su acuerdo con la generación voraz $g$ mediante la \textbf{similitud de Jaccard
de $n$-gramas de carácter}:
\begin{equation}
\mathrm{sim}(a,b) = \frac{|N(a)\cap N(b)|}{|N(a)\cup N(b)|},\qquad
c = \frac{1}{K}\sum_{k=1}^{K}\mathrm{sim}(g,s_k),
\end{equation}
donde $N(\cdot)$ es el conjunto de $4$-gramas de carácter. Una baja consistencia $c$ indica un
contenido inestable entre muestras, señal candidata de alucinación.

\paragraph{Dependencia de la imagen.} Se genera además un informe con la \textbf{imagen en blanco}
$g_{\varnothing}$ y se calcula $\mathrm{blank\_sim} = \mathrm{sim}(g, g_{\varnothing})$: cuanto menor
es, más depende de la imagen la salida real.

\paragraph{Anclaje objetivo.} Sobre NIH, dado el mapa de calor normalizado $h$ y la caja $B$ de área
relativa $|B|$, el \textbf{lift} mide la mejora sobre el azar de la energía de atención contenida en
la caja:
\begin{equation}
\mathrm{lift} = \frac{\sum_{(i,j)\in B} h_{ij}}{|B|}.
\end{equation}
Un valor $\mathrm{lift}>1$ indica que la atención se concentra en la lesión por encima de lo
esperable por azar.

\section{Protocolo de etiquetado del error}
La variable dependiente —\emph{correcto}, \emph{alucinación} u \emph{omisión}— se determina por dos
mecanismos complementarios. El \textbf{etiquetado objetivo} compara los hallazgos agudos afirmados por
el modelo con la verdad estructurada (MeSH en Open-I; presentes/ausentes en PadChest-GR): se cuenta
alucinación cuando el modelo afirma un hallazgo agudo ausente de la referencia, y omisión cuando
omite uno presente. El \textbf{etiquetado por juez} (G-Eval) somete cada par (informe generado,
referencia) a un modelo de lenguaje que aplica una \textbf{rúbrica fija}: se consideran significativos
únicamente los hallazgos \emph{agudos} (cardiomegalia, derrame, consolidación, edema, neumotórax,
masa, nódulo, neumonía, infiltrado, atelectasia), ignorándose los crónicos, óseos, vasculares,
técnicos o triviales; un informe es \emph{correcto} si coincide con la referencia en los hallazgos
significativos, \emph{alucinación} si afirma uno ausente y \emph{omisión} si pierde uno presente. La
concordancia entre ambos mecanismos valida el etiquetado y mitiga la circularidad de depender de un
único criterio.

\section{Requisitos del sistema}
Aunque el trabajo es fundamentalmente de investigación, el sistema resultante se especifica mediante
requisitos, lo que ordena su diseño y su evaluación. Los \textbf{requisitos funcionales} (RF) son:
\textbf{RF1}, generar el informe de una radiografía de tórax; \textbf{RF2}, extraer en inferencia
las señales de incertidumbre sin generaciones adicionales obligatorias; \textbf{RF3}, calcular el
anclaje de la explicación (mapa de calor o cajas); \textbf{RF4}, estimar y \textbf{calibrar} una
confianza por informe; \textbf{RF5}, presentar informe, anclaje y confianza de forma integrada; y
\textbf{RF6}, priorizar la revisión mediante bandas de riesgo. Los \textbf{requisitos no funcionales}
(RNF) son: \textbf{RNF1}, ejecutarse en una única tarjeta de 12\,GB mediante cuantización;
\textbf{RNF2}, que la señal de confianza sea \textbf{agnóstica al modelo} y no incurra en coste de
inferencia adicional; \textbf{RNF3}, reproducibilidad (extracción resumible, cálculo determinista);
\textbf{RNF4}, interpretabilidad y honestidad de la salida (comunicar las limitaciones); y
\textbf{RNF5}, portabilidad y bajo coste de despliegue.

\section{Arquitectura de la solución}
La solución se organiza en \textbf{cinco capas}. La \textbf{capa de datos} gestiona los conjuntos
(Open-I, NIH, PadChest-GR), su descarga, anonimización y preprocesado. La \textbf{capa de modelo e
inferencia} carga los generadores cuantizados y produce, por imagen, el informe, las probabilidades y
entropías por token, las muestras estocásticas y la atención visual. La \textbf{capa de análisis y
explicabilidad} calcula a partir de lo anterior las señales de incertidumbre, las métricas de anclaje
y las etiquetas de error. La \textbf{capa de calibración} transforma la señal en una probabilidad con
bandas de riesgo (el semáforo). Y la \textbf{capa de presentación} materializa el sistema en la
demostración interactiva. Esta separación por capas favorece la reutilización —el marco de análisis es
\emph{agnóstico al generador}— y facilita sustituir un modelo por otro sin alterar el resto. La
Figura~\ref{fig:pipeline} representa el flujo completo del sistema.

\begin{figure}[t]\centering
\begin{tikzpicture}[
  node distance=6mm and 9mm,
  box/.style={draw, rounded corners, align=center, minimum height=9mm, minimum width=20mm, font=\footnotesize, fill=blue!5},
  data/.style={draw, align=center, minimum height=9mm, font=\footnotesize, fill=gray!8},
  outbox/.style={draw, rounded corners, align=center, minimum height=9mm, font=\footnotesize, fill=green!8},
  arr/.style={-{Latex[length=2mm]}, thick}]
  \node[data] (img) {Radiografía};
  \node[box, right=of img] (gen) {Generador\\(LLaVA-Rad /\\MAIRA-2)};
  \node[data, right=of gen] (sig) {Informe\\Probabilidades\\Atención};
  \node[box, right=of sig] (an) {Análisis:\\señales +\\anclaje};
  \node[box, right=of an] (cal) {Calibración\\(semáforo)};
  \node[outbox, below=8mm of cal] (demo) {Demo: informe\\+ anclaje\\+ confianza};
  \draw[arr] (img) -- (gen);
  \draw[arr] (gen) -- (sig);
  \draw[arr] (sig) -- (an);
  \draw[arr] (an) -- (cal);
  \draw[arr] (cal) -- (demo);
  \node[data, below=8mm of an] (truth) {Verdad:\\MeSH / radiólogo\\/ cajas NIH};
  \draw[arr] (truth) -- (an);
\end{tikzpicture}
\caption{Flujo del sistema: de la radiografía al informe con anclaje y confianza calibrada. El bloque
de análisis se valida contra la verdad objetiva (MeSH, hallazgos de radiólogo, cajas de NIH).}
\label{fig:pipeline}
\end{figure}

\section{Casos de uso}
Se contemplan cuatro casos de uso principales. \textbf{CU-01 (radiólogo): obtener informe y confianza}
—el usuario carga una radiografía y el sistema devuelve el informe, el anclaje y la banda de
confianza—. \textbf{CU-02 (radiólogo): consultar el anclaje} —inspeccionar dónde sustenta el modelo su
afirmación, como apoyo a la lectura—. \textbf{CU-03 (radiólogo): priorizar la revisión} —ordenar la
carga de trabajo atendiendo primero a los informes en zona roja—. Y \textbf{CU-04 (investigador):
auditar y reproducir} —extraer las señales sobre un conjunto y reproducir las métricas de anclaje y
detección—.

\section{Planificación y desarrollo}
El proyecto se ha desarrollado de forma incremental, con control de versiones (Git) y gestión de
tareas por incrementos, organizándose en seis fases que se describen a continuación. Cada fase produce
un artefacto verificable —un componente de software o un resultado experimental— antes de avanzar a la
siguiente, lo que permite congelar resultados y retomar procesos largos entre sesiones. El trabajo de
campo experimental (extracción de señales y etiquetado sobre miles de imágenes) se ejecuta en procesos
reanudables que sobreviven a interrupciones, dado el tiempo de cómputo en hardware de consumo.

\section{Diseño experimental}
El trabajo se estructura en seis fases experimentales.
\begin{enumerate}
  \item \textbf{Modelo propio (POC).} Un codificador visual (Vision Transformer) y un decodificador
  de lenguaje (GPT-2) al que se inyecta atención cruzada en cada bloque mediante una compuerta
  inicializada a cero, que preserva los pesos preentrenados. El entrenamiento sigue dos etapas:
  alineamiento contrastivo (InfoNCE) y ajuste generativo. Se realiza una ablación (M0--M3) que
  introduce progresivamente la destilación de conocimiento textual (CXR-BERT
  \parencite{boecking2022}) y visual (RAD-DINO \parencite{perezgarcia2024}).
  \item \textbf{Elevación a un modelo del estado del arte.} El marco de explicabilidad se aplica a
  LLaVA-Rad, ejecutado en el hardware local mediante cuantización a 8 bits, adaptando la
  extracción de mapas a su arquitectura de tokens visuales antepuestos.
  \item \textbf{Anclaje objetivo.} Validación del anclaje contra las cajas de NIH.
  \item \textbf{Detección de errores.} Estudio multi-señal de la fiabilidad, combinando morfología,
  incertidumbre, dependencia de imagen y consistencia, con etiqueta objetiva (MeSH) y juez.
  \item \textbf{Réplica cruzada (diseño \(2\times2\)).} La señal se valida sobre un \textbf{segundo
  generador del estado del arte, MAIRA-2} \parencite{maira2}, y sobre un \textbf{segundo dominio,
  PadChest-GR}, cruzando ambos factores (dos modelos \(\times\) dos dominios). MAIRA-2 se ejecuta
  también cuantizado a 8 bits; al haber sido entrenado con PadChest, su evaluación
  \emph{in-distribution} se realiza sobre PadChest y la evaluación \emph{fuera de distribución} sobre
  Open-I, ajeno a su entrenamiento. Se aprovecha además su capacidad de \textbf{anclaje explícito
  (cajas)} para contrastar por tercera vía si la localización de la explicación predice la fiabilidad.
  \item \textbf{Calibración y confianza operativa.} La señal final se calibra para expresar la
  fiabilidad como una probabilidad con bandas de riesgo (``semáforo''), y se materializa en una
  \textbf{demostración interactiva} que integra informe, anclaje y confianza.
\end{enumerate}

\section{Métricas de evaluación}
La evaluación cuantitativa se apoya en un conjunto de métricas que conviene definir formalmente.

\paragraph{Discriminación (AUC--ROC).} La capacidad de una señal para separar dos clases se mide con el
\textbf{área bajo la curva ROC}. La curva ROC representa la sensibilidad (tasa de verdaderos positivos)
frente a la tasa de falsos positivos al variar el umbral de decisión. El AUC admite una interpretación
probabilística: es la probabilidad de que una instancia positiva reciba una puntuación mayor que una
negativa elegida al azar, y equivale al estadístico de Mann--Whitney,
\begin{equation}
\mathrm{AUC} = \frac{1}{n_{+}n_{-}}\sum_{i\in\mathcal{P}}\sum_{j\in\mathcal{N}} \mathbf{1}[s_i > s_j],
\end{equation}
donde $\mathcal{P}$ y $\mathcal{N}$ son los conjuntos de positivos y negativos y $s$ las puntuaciones.
Un valor de $0{,}5$ indica ausencia de discriminación y $1{,}0$ discriminación perfecta. Se calcula
mediante rangos para garantizar reproducibilidad.

\paragraph{Significancia (test de permutación).} Para contrastar si un AUC observado supera al azar
—especialmente con clases desbalanceadas— se emplea un \textbf{test de permutación}: se permutan
aleatoriamente las etiquetas un número elevado de veces, se recalcula el AUC en cada permutación y se
obtiene el $p$-valor como la fracción de valores nulos que igualan o superan al observado.

\paragraph{Intervalos de confianza (\textit{bootstrap}).} La incertidumbre de la estimación se cuantifica
por \textbf{\textit{bootstrap} estratificado}: se remuestrean con reemplazo, por separado, los
positivos y los negativos un número elevado de veces (2000 en este trabajo), y el intervalo de
confianza al 95\,\% se toma de los percentiles $2{,}5$ y $97{,}5$ de la distribución de AUC
remuestreados.

\paragraph{Calibración.} Que una señal discrimine no implica que esté \textbf{calibrada}. La
calibración se evalúa con la \textbf{curva de fiabilidad}, que compara la probabilidad predicha con la
tasa de acierto observada por grupos, y con el \textbf{\textit{Brier score}}, el error cuadrático medio
entre la probabilidad predicha $\hat{p}_i$ y el resultado $y_i$:
$\mathrm{Brier}=\frac{1}{n}\sum_i (\hat{p}_i - y_i)^2$. Sobre esta base se definen las bandas de riesgo
del semáforo.

\paragraph{Anclaje y similitud.} El anclaje se mide con el \textit{lift}, el \textit{energy-in-bbox} y
el \textit{pointing-game} ya definidos; la consistencia textual, con la similitud de Jaccard de
$n$-gramas de carácter.

\section{Método científico y control de confusores}
El trabajo sigue un enfoque \textbf{hipotético-deductivo}: se formulan hipótesis falsables (H1--H4,
Sección~4.2) y se diseñan experimentos que puedan refutarlas. Tres principios metodológicos guían la
investigación. Primero, la \textbf{validación contra verdad objetiva}: siempre que es posible, las
señales se contrastan con referencias externas (cajas reales, diagnóstico estructurado) en lugar de con
otras señales del propio modelo, y el etiquetado del error se \textbf{triangula} con tres fuentes
independientes. Segundo, el \textbf{control de confusores}: se identifican y neutralizan las variables
espurias que podrían simular una señal —la longitud del informe (por residualización), la posición del
token (usando la probabilidad de contenido) o el tipo de mapa de calor (mediante comparaciones
emparejadas)—. Tercero, la \textbf{replicación}: los hallazgos se comprueban en un diseño cruzado de dos
modelos por dos dominios antes de darlos por robustos. Este rigor resultó decisivo: varias señales de
apariencia prometedora (el máximo de atención, la entropía cruda, la relación \emph{lift}--acierto) se
revelaron \textbf{artefactos} que solo un control cuidadoso de confusores permitió descartar,
subrayando la importancia del método sobre el resultado aparente.

\section{Modelo de análisis de datos}
El análisis es fundamentalmente cuantitativo. La capacidad discriminativa de cada señal se mide
mediante AUC univariante; las señales se combinan con regresión logística evaluada por
\textbf{validación cruzada out-of-fold} —de modo que cada predicción se obtiene sobre datos no
empleados en el ajuste— y su significancia se contrasta con \textbf{tests de permutación} que
controlan el desbalance de clases. Se prestó especial atención al control de \textit{confounds}
(régimen de la muestra, longitud del informe y artefactos del máximo de atención), validando cada
señal frente a una referencia continua u objetiva para descartar espejismos. La exploración de
señales adicionales se realiza \textbf{sobre los tres conjuntos a la vez}, exigiendo que cualquier
mejora \emph{generalice}; se comprueba que ninguna combinación multi-señal supera de forma
consistente a la probabilidad mínima de contenido por sí sola —las demás señales resultan específicas
de modelo o de dominio—, por lo que la fórmula final privilegia la \textbf{simplicidad robusta} de
una única señal frente a un combinado sobreajustado. Esta se \textbf{calibra} para expresar la
fiabilidad como una probabilidad con su \textbf{intervalo de confianza (\emph{bootstrap})} y bandas
de riesgo. El software empleado es Python (PyTorch, \texttt{transformers}, NumPy) y la cuantización se
realiza con \texttt{bitsandbytes}; el AUC y los intervalos se calculan con estimadores basados en
rangos para garantizar reproducibilidad.

\chapter{Presentación de análisis, resultados y discusión}
% Redacción con números finales (2026-08-10): 2x2 modelo x dominio, MAIRA-2, token de apertura,
% calibración y demos. Estilo impersonal UAX. Cada figura se justifica explícitamente en el texto.

En este capítulo se describe la muestra analizada, se presentan los resultados y se discute su
significado. La línea argumental es la siguiente: la atención del modelo está clínicamente anclada,
pero su morfología no permite estimar la fiabilidad; esta reside en la \emph{incertidumbre de la
decisión} del propio modelo y se detecta, sobre todo, para un tipo de error —las alucinaciones—. Se
demuestra que la señal es \emph{agnóstica al modelo} y \emph{robusta al dominio} validándola sobre
dos generadores (LLaVA-Rad y MAIRA-2) y dos dominios (Estados Unidos y España), y se traduce en una
confianza calibrada apta para el cribado clínico.

\section{Análisis descriptivo de la muestra}
El estudio de fiabilidad se articula como un diseño \(2\times2\) que cruza \textbf{dos generadores}
—LLaVA-Rad (7B, entrenado sobre MIMIC-CXR) y MAIRA-2 (7B, entrenado sobre MIMIC-CXR, PadChest y
USMix)— con \textbf{dos dominios}: Open-I (Indiana University, EE.\,UU.) y PadChest-GR (España). Para
cada celda se generan informes y se etiqueta cada uno como \emph{correcto}, \emph{alucinación}
(afirmar un hallazgo agudo ausente) u \emph{omisión} o \emph{miss} (perder un hallazgo agudo
presente). El etiquetado combina tres fuentes de verdad independientes para evitar la circularidad
de un único juez: diagnóstico estructurado oficial (MeSH de Open-I; hallazgos presentes/ausentes de
PadChest-GR) y juicio basado en modelo de lenguaje sobre la narrativa completa (\emph{G-Eval}). El
anclaje de la atención se evalúa aparte sobre las \textbf{984 anotaciones con caja} de NIH
ChestX-ray8. La Tabla~\ref{tab:muestra} resume las cuatro celdas. La prevalencia de error y el
predominio de informes normales son característicos del dominio; la cardiomegalia sobre-diagnosticada
en tórax normal aparece como sesgo sistemático de \emph{ambos} generadores.

\begin{table}[t]\centering
\caption{Muestra del diseño \(2\times2\) (etiquetas objetivas/G-Eval). n = informes; alu = alucinaciones; mis = omisiones.}
\label{tab:muestra}
\begin{tabular}{llrrr}
\toprule
Generador & Dominio & n & alu & mis \\ \midrule
LLaVA-Rad & Open-I (in-distribution)  & 800 & 110 & 228 \\
LLaVA-Rad & PadChest (OOD)            & 381 & 33  & 130 \\
MAIRA-2   & PadChest (in-distribution)& 379 & 27  & 80  \\
MAIRA-2   & Open-I (OOD)              & 300 & 12  & 22  \\ \bottomrule
\end{tabular}
\end{table}

\section{El generador propio y la lección sobre las métricas}
El generador propio se desarrolló de forma deliberada como paso previo a emplear un modelo del
estado del arte ya entrenado: construirlo desde cero permite observar su \textbf{funcionamiento
interno} —cómo se integra la señal visual en la generación de texto, y dónde falla— de un modo que
un modelo de caja cerrada no ofrece, y expone de primera mano las \textbf{limitaciones propias de un
modelo de menor capacidad} entrenado con datos escasos. En su implementación se documentaron y
resolvieron dos fallos característicos. El primero, el \textbf{colapso de la generación} al inyectar la atención
cruzada: la contribución visual, inicialmente ruidosa, desestabilizaba el decodificador preentrenado;
se mitigó con la \textbf{compuerta inicializada a cero}, que preserva el comportamiento del modelo de
lenguaje y deja que la señal visual se incorpore de forma gradual. El segundo, la \textbf{ceguera a la
imagen}: el decodificador tendía a ignorar la imagen y a reproducir la plantilla mayoritaria; se
mitigó congelando el decodificador base y entrenando solo la atención cruzada (PEFT).

Sobre esta base se diseñó una ablación (Tabla~\ref{tab:ablacion}) que introduce progresivamente la
destilación de conocimiento textual y visual, con el objetivo de medir su efecto sobre la calidad y el
anclaje. La lección metodológica más relevante emergió al evaluar estas variantes: la métrica
\textbf{ROUGE engaña}. Una variante mejor anclada a la imagen —que describe hallazgos reales— puede
obtener un ROUGE \emph{inferior} al de un modelo ciego que memoriza la plantilla ``normal'', porque en
un dominio dominado por informes normales el solapamiento léxico premia la reproducción de la
plantilla, no la lectura de la imagen. Este hallazgo, coherente con la crítica del Capítulo~3, motivó
el abandono de las métricas de solapamiento en favor de la \textbf{corrección clínica} y, en última
instancia, el giro de todo el trabajo hacia la \textbf{fiabilidad}. Junto con las limitaciones de
capacidad ya observadas en el propio modelo, esta lección sobre las métricas es la que llevó a
\textbf{elevar el análisis a un MLLM del estado del arte} (LLaVA-Rad): solo un generador de calidad
suficiente permite separar los errores reales de los artefactos de un modelo de escala reducida, y
solo evaluando la corrección clínica —no el solapamiento léxico— tiene sentido preguntarse si la
explicación o la incertidumbre predicen la fiabilidad.

\begin{table}[t]\centering
\caption{Ablación de destilación del generador propio (M0--M3): fuentes de conocimiento introducidas.}
\label{tab:ablacion}
\begin{tabular}{lll}
\toprule
Variante & Destilación textual & Destilación visual \\ \midrule
M0 (línea base) & --- & --- \\
M1 & CXR-BERT (torre congelada) & --- \\
M2 & CXR-BERT (destilación al \emph{student}) & --- \\
M3 & CXR-BERT (destilación) & RAD-DINO (CLS $\to$ CLS) \\ \bottomrule
\end{tabular}
\end{table}

\section{La atención está clínicamente anclada}
Contrastada con las cajas reales de NIH, la atención de LLaVA-Rad se concentra sobre la lesión muy
por encima del azar: \textbf{lift global 3,03} y \emph{pointing-game} 0,38 (3,7 veces el azar). El
anclaje varía con sentido clínico (Tabla~\ref{tab:nih}): fuerte en hallazgos focales pequeños
(nódulo, masa), débil en la cardiomegalia, difusa y de caja grande. Este resultado justifica el uso
del mapa de calor como puntero de \emph{dónde} mira el modelo, pero —como se verá— no como medida de
\emph{cuánto} fiarse.

\begin{table}[t]\centering
\caption{Anclaje objetivo sobre NIH (n = 388 hallazgos detectados). Lift = energía-en-caja / azar.}
\label{tab:nih}
\begin{tabular}{lrrr}
\toprule
Patología & Detección & Lift & Pointing \\ \midrule
Nódulo & 0,22 & 7,91 & 0,47 \\ Masa & 0,18 & 5,10 & 0,87 \\
Atelectasia & 0,29 & 3,45 & 0,30 \\ Derrame & 0,58 & 3,23 & 0,36 \\
Neumotórax & 0,74 & 2,81 & 0,37 \\ Neumonía & 0,29 & 2,67 & 0,23 \\
Infiltrado & 0,19 & 2,62 & 0,48 \\ Cardiomegalia & 0,58 & 1,63 & 0,39 \\ \midrule
\textbf{Global} & 0,39 & \textbf{3,03} & 0,38 \\ \bottomrule
\end{tabular}
\end{table}

La Figura~\ref{fig:nih} ilustra visualmente este anclaje: para cada patología se muestran los casos de
mayor \emph{lift}, con el mapa de calor superpuesto y la caja real. Se observa que la atención se
concentra sobre la lesión, de forma especialmente nítida en los hallazgos focales pequeños (nódulo,
masa, neumotórax) y más difusa en la cardiomegalia, de caja grande.

\begin{figure}[p]\centering
\includegraphics[height=0.92\textheight]{fig_nih_montage.png}
\caption{Anclaje de la atención sobre NIH por patología (casos de mayor \emph{lift}). Mapa de calor
superpuesto y caja real (verde); se indica el \emph{lift}, el acierto/fallo del veredicto y el área
relativa de la caja. La atención se concentra sobre la lesión, sobre todo en hallazgos focales.}
\label{fig:nih}
\end{figure}

A modo de contraste, la Figura~\ref{fig:nihmed} muestra los casos de \emph{lift} mediano por patología:
el anclaje sigue siendo apreciable pero más difuso, ilustrando que la calidad de la localización varía
de forma continua y que no todos los aciertos van acompañados de una atención igualmente concentrada.

\begin{figure}[p]\centering
\includegraphics[height=0.92\textheight]{fig_nih_median.png}
\caption{Anclaje sobre NIH por patología en casos de \emph{lift} mediano. La atención sigue sobre la
región de la lesión pero de forma más difusa que en los casos de la Figura~\ref{fig:nih}.}
\label{fig:nihmed}
\end{figure}

\section{La morfología de la explicación no predice la fiabilidad}
Ninguna característica geométrica del mapa de calor (concentración, dispersión, agrupamiento) separa
de forma robusta los informes correctos de los erróneos (AUC 0,54; no significativo). Las señales
que aparentaban serlo resultaron artefactos detectados al validarlas: la aparente relación entre
anclaje (\emph{lift}) y acierto \textbf{desaparece} (AUC 0,52) al emplear un mapa de calor homogéneo,
revelándose como un artefacto de contrastar mapas por-hallazgo con mapas globales. El resultado se
confirma por tercera vía con las \textbf{cajas de anclaje propias de MAIRA-2}: el modelo localiza sus
hallazgos \emph{inventados} con cajas anatómicamente plausibles (una caja sobre el corazón para una
cardiomegalia inexistente, sobre el pulmón para una enfermedad intersticial inexistente), de modo que
la presencia y la forma del anclaje acompañan a la \emph{afirmación}, no al \emph{acierto} (frecuencia
de caja del 19\,\% en alucinaciones frente al 7\,\% en informes correctos, dirección inútil como
señal). En resumen, ni el mapa de calor de atención, ni el \emph{lift} sobre NIH, ni las cajas
entrenadas del propio modelo permiten leer la fiabilidad en la morfología de la explicación.

\section{La incertidumbre de la decisión detecta las alucinaciones}
Frente a la morfología, la \textbf{incertidumbre generativa} —operacionalizada como la probabilidad
mínima de \emph{contenido} de los tokens del informe— sí discrimina la alucinación del acierto. La
Figura~\ref{fig:sep} muestra, para las cuatro celdas, las distribuciones de esta señal en informes
correctos, con omisión y con alucinación. \emph{Se incluye esta figura porque hace visible el
resultado central}: las alucinaciones (rojo) se desplazan hacia la baja confianza y se separan de los
informes correctos (verde), con la única excepción de LLaVA-Rad fuera de distribución, donde ambas
distribuciones se solapan.

\begin{figure}[t]\centering
\includegraphics[width=\linewidth]{fig_separacion_clases.png}
\caption{Separación de clases de la probabilidad mínima de contenido (menor $=$ más incierto) en las
cuatro celdas del diseño. Verde: correcto; gris: omisión; rojo: alucinación; líneas discontinuas:
medias. Las alucinaciones se concentran en la zona incierta salvo en LLaVA-Rad$\cdot$PadChest (OOD).}
\label{fig:sep}
\end{figure}

\paragraph{El token de apertura contamina la señal (aportación metodológica).}
La probabilidad mínima calculada sobre \emph{todo} el informe rinde de forma mediocre porque el
mínimo recae casi siempre en el token de \emph{apertura} del informe (``The'', ``Frontal''\dots),
que refleja incertidumbre de estilo, no de contenido, y es uniformemente baja en todos los casos.
Excluir ese primer token —usar la probabilidad mínima de contenido— eleva sustancialmente la
detección en las cuatro celdas (Figura~\ref{fig:token}). \emph{Se justifica esta figura porque aísla
una decisión de diseño con impacto grande y generalizable}: el AUC pasa de 0,71 a 0,90 en Open-I y de
0,50 a 0,80 en MAIRA-2$\cdot$PadChest, sin coste alguno. La interpretación de este efecto es
instructiva: el primer token del informe concentra la incertidumbre \emph{estilística} —qué fórmula de
apertura emplear (``The'', ``Frontal''\dots)—, que es alta y uniforme en todos los informes con
independencia de su contenido, y que por tanto \emph{no discrimina} y solo añade ruido al mínimo. Al
restringir la señal al contenido, esta pasa a reflejar la incertidumbre \emph{diagnóstica} del modelo.
Más allá de la mejora numérica, el hallazgo ilustra un principio general aplicable a otras tareas de
generación: al resumir la incertidumbre de una secuencia conviene separar la duda sobre la forma de la
duda sobre el fondo, pues solo la segunda informa sobre la fiabilidad del contenido.

\begin{figure}[t]\centering
\includegraphics[width=0.92\linewidth]{fig_token_apertura.png}
\caption{Efecto de excluir el token de apertura. Barras grises: probabilidad mínima sobre todo el
informe; barras verdes: sobre el contenido (sin el primer token). La mejora es sistemática en las
cuatro celdas.}
\label{fig:token}
\end{figure}

\paragraph{La señal es agnóstica al modelo y robusta al dominio.}
La Figura~\ref{fig:2x2} sintetiza el resultado principal como matriz \(2\times2\). \emph{Se incluye
porque resume, en una sola imagen, la validez externa del método}: la detección es fuerte in-distribution
(LLaVA-Rad 0,90; MAIRA-2 0,80) y, crucialmente, \textbf{se mantiene fuera de distribución para el
modelo de mayor calidad} (MAIRA-2 0,81, prácticamente igual que su valor in-distribution). La única
celda débil es LLaVA-Rad fuera de distribución (0,61): un generador limitado, al enfrentarse a
imágenes ajenas a su entrenamiento, recae en un informe-plantilla emitido con confianza, lo que
erosiona la señal. La degradación fuera de distribución es, por tanto, una limitación del generador
débil, no del método. Este matiz es importante porque invierte la conclusión que se habría extraído de
haber estudiado un único generador: con LLaVA-Rad aislado, la caída a 0,61 fuera de distribución
sugeriría que la fiabilidad por incertidumbre no generaliza; el diseño cruzado revela que sí lo hace
cuando el modelo es lo bastante bueno, y que el problema reside en que un generador débil, ante
imágenes ajenas a su entrenamiento, deja de leer la imagen y ``rellena'' con su plantilla más
probable, emitida con confianza. La lección práctica es doble: conviene, primero, elegir un generador
de calidad, y segundo, calibrar y verificar la señal en cada dominio de despliegue antes de confiar en
ella.

\begin{figure}[t]\centering
\includegraphics[width=0.72\linewidth]{fig_2x2_auc.png}
\caption{AUC de detección de alucinaciones (probabilidad mínima de contenido) en el diseño
\(2\times2\). Verde $=$ alto. La señal es robusta salvo en el generador débil fuera de distribución.}
\label{fig:2x2}
\end{figure}

Las curvas ROC de la Figura~\ref{fig:roc} caracterizan el compromiso sensibilidad--especificidad de
la señal en las cuatro celdas. \emph{Se muestran porque traducen el AUC agregado en el comportamiento
operativo}: las tres celdas robustas dominan claramente la diagonal de azar, mientras que
LLaVA-Rad$\cdot$PadChest se aproxima a ella, coherente con el solapamiento de la Figura~\ref{fig:sep}.

\begin{figure}[t]\centering
\includegraphics[width=0.62\linewidth]{fig_roc.png}
\caption{Curvas ROC de detección de alucinaciones (probabilidad mínima de contenido). Tres de las
cuatro celdas dominan la diagonal de azar; LLaVA-Rad$\cdot$PadChest (OOD) se aproxima a ella.}
\label{fig:roc}
\end{figure}

\subsection{Ninguna señal adicional generaliza: en favor de la simplicidad}
Se exploró si combinar la probabilidad mínima de contenido con otras señales agnósticas (entropía,
verosimilitud media, auto-consistencia, dependencia de imagen) mejoraba la detección. La exploración
se realizó exigiendo que cualquier mejora \emph{generalizara} sobre los tres conjuntos con suficiente
número de alucinaciones. La Tabla~\ref{tab:senales} muestra el AUC univariante de cada señal por
conjunto: ninguna supera de forma consistente a la probabilidad mínima de contenido. Señales fuertes
en los conjuntos de LLaVA-Rad (entropía, verosimilitud) resultan inútiles en MAIRA-2, y la dependencia
de imagen ayuda en MAIRA-2 pero se invierte en Open-I; son, por tanto, específicas de modelo o de
dominio. En consecuencia, la fórmula final privilegia la \textbf{simplicidad robusta} de una única
señal frente a un combinado sobreajustado que no transferiría.

\begin{table}[t]\centering
\caption{AUC univariante de cada señal (orientada: mayor $=$ más probable alucinación) por conjunto.
Ninguna mejora de forma consistente a la probabilidad mínima de contenido.}
\label{tab:senales}
\begin{tabular}{lrrr}
\toprule
Señal & Open-I (LLaVA) & PadChest (LLaVA) & PadChest (MAIRA) \\ \midrule
\textbf{min\_prob de contenido} & \textbf{0,90} & \textbf{0,57} & \textbf{0,80} \\
Percentil 5 de probabilidad & 0,90 & 0,61 & 0,58 \\
Verosimilitud media & 0,89 & 0,65 & 0,49 \\
Entropía media & 0,92 & 0,65 & 0,46 \\
Auto-consistencia & 0,73 & 0,43 & 0,65 \\
Dependencia de imagen (blank) & 0,32 & 0,51 & 0,70 \\ \bottomrule
\end{tabular}
\end{table}

\section{Disociación alucinación--omisión}
La detección es \emph{asimétrica}. La Figura~\ref{fig:diso} contrasta, por celda, el AUC para separar
alucinaciones frente a correctos y omisiones frente a correctos. \emph{Se justifica esta figura porque
delimita con honestidad el alcance del método}: la señal detecta las alucinaciones de forma robusta
(0,80--0,90), pero las omisiones solo de forma parcial y variable (0,47--0,85). Las omisiones no son
siempre errores ``confiados''; como se aprecia en la Figura~\ref{fig:sep} (\mbox{distribución} gris), tienden
a situarse \emph{entre} los correctos y las alucinaciones, cerca del umbral, y fuera de distribución
(MAIRA-2$\cdot$Open-I) se desplazan hacia la incertidumbre. La interpretación es que, cuando el modelo
fabrica un hallazgo, titubea; cuando lo omite emitiendo un informe normal, a menudo lo hace con
confianza próxima a la de un informe normal correcto, aunque no siempre. En consecuencia, la fórmula
es un \emph{detector de alucinaciones} que captura además una fracción de omisiones como beneficio
colateral, no un detector general de error.

\begin{figure}[t]\centering
\includegraphics[width=0.82\linewidth]{fig_disociacion.png}
\caption{Disociación: AUC para separar alucinaciones (rojo) y omisiones (gris) de los informes
correctos, por celda. La incertidumbre detecta alucinaciones de forma robusta y omisiones solo
parcialmente.}
\label{fig:diso}
\end{figure}

\section{Calibración y confianza operativa}
Para convertir la señal en una salida clínica se estudia su calibración. La Figura~\ref{fig:cal}
representa la tasa de alucinación observada en función de la probabilidad mínima de contenido:
\emph{se incluye porque es el fundamento empírico del ``semáforo'' de confianza}, al mostrar una
relación monótona —a menor confianza del modelo, mayor tasa de alucinación observada— sobre la que se
definen bandas de riesgo. Sobre Open-I, las bandas por cuartiles del grupo correcto arrojan una tasa
de alucinación del 30\,\% en la banda roja (revisar), 6\,\% en la amarilla y $\approx$0\,\% en las
verdes. Como priorizador de revisión (Tabla~\ref{tab:oper}), un umbral de Youden captura el 84\,\% de
las alucinaciones —y, además, el 44\,\% de las omisiones, como beneficio colateral— revisando en
torno a un tercio de los estudios; a especificidad del 95\,\%, el 53\,\% de las alucinaciones revisando solo el 5\,\%. La
propuesta no es, pues, un clasificador binario —la baja prevalencia lo haría poco preciso— sino una
\textbf{confianza calibrada con intervalo} que ordena la carga de revisión, con el radiólogo como
decisor.

\begin{figure}[t]\centering
\includegraphics[width=0.66\linewidth]{fig_calibracion.png}
\caption{Calibración: tasa de alucinación observada frente a la probabilidad mínima de contenido
(media por decil, datos agregados). La relación monótona fundamenta las bandas del semáforo de
confianza.}
\label{fig:cal}
\end{figure}

\begin{table}[t]\centering
\caption{Puntos de operación sobre Open-I (verdad MeSH; puntuación de contenido). FA $=$ falsas alarmas.}
\label{tab:oper}
\begin{tabular}{lrrr}
\toprule
Umbral & Recall alucinación & Recall omisión & FA \\ \midrule
Youden          & 84\,\% & 44\,\% & 15\,\% \\
Especificidad 90\,\% & 71\,\% & 33\,\% & 10\,\% \\
Especificidad 95\,\% & 53\,\% & 20\,\% & 5\,\%  \\ \bottomrule
\end{tabular}
\end{table}

\section{Traducción clínica: la confianza como semáforo}
Los resultados anteriores solo tienen valor si pueden \emph{presentarse} al clínico de una forma que
mejore su trabajo sin introducir nuevos riesgos. Esta sección justifica las decisiones de diseño de la
demostración desarrollada y su encaje en un entorno médico real.

\subsection{Por qué una demostración interactiva}
Se ha optado por materializar el sistema en una \textbf{demostración interactiva} (con la biblioteca
Gradio) que, ante una radiografía, muestra en un único panel tres elementos: el \textbf{informe}
generado, el \textbf{anclaje} (dónde mira el modelo, mediante mapa de calor o cajas) y la
\textbf{confianza} calibrada (cuánto fiarse). La elección responde a tres criterios. Primero, la
\textbf{reproducibilidad y accesibilidad}: una interfaz web ligera permite que un radiólogo,
un investigador o un evaluador reproduzcan el comportamiento del sistema sin conocimientos técnicos ni
instalación compleja. Segundo, la \textbf{integración de las tres señales} en una sola vista, que es
precisamente la propuesta del trabajo: informe, evidencia visual y fiabilidad no deben consultarse por
separado. Y tercero, su idoneidad como \textbf{prototipo para un estudio de lectura} futuro, al
permitir presentar casos controlados y recoger la interacción del clínico.

\subsection{Un gradiente de confianza, no una alerta de error}
La decisión de diseño más relevante es presentar la fiabilidad como un \textbf{semáforo en gradiente}
—una banda continua de verde a rojo— en lugar de una \textbf{alerta binaria} de posible error. La
justificación es doble. Por un lado, un clasificador binario sería \emph{poco preciso} dada la baja
prevalencia de las alucinaciones: generaría muchas falsas alarmas por cada acierto (véase
Tabla~\ref{tab:oper}). Por otro, y más importante, la literatura de sistemas de ayuda a la decisión
clínica documenta el fenómeno de la \textbf{fatiga de alarmas}: cuando un sistema emite alertas
frecuentes —muchas de ellas irrelevantes— los profesionales se \emph{desensibilizan} y terminan
ignorándolas, anulando su utilidad e incluso introduciendo riesgo. Un gradiente calibrado evita este
efecto: la mayoría de los informes fiables permanecen en verde, de baja fricción, y la atención del
radiólogo se \textbf{prioriza} hacia la minoría de casos de riesgo elevado, en lugar de dispersarse en
una sucesión de alertas. La confianza se convierte así en un \emph{ordenador de la carga de revisión},
coherente con los puntos de operación de la Tabla~\ref{tab:oper}. La Figura~\ref{fig:semaforo}
representa el semáforo con las bandas calibradas sobre Open-I.

\begin{figure}[t]\centering
\includegraphics[width=0.9\linewidth]{fig_semaforo.png}
\caption{Semáforo de confianza en gradiente. Las bandas, definidas sobre los cuartiles de la
distribución de los informes correctos, se anotan con la probabilidad de alucinación observada. El
verde concentra los informes fiables (baja fricción) y el rojo prioriza la revisión.}
\label{fig:semaforo}
\end{figure}

La Figura~\ref{fig:demo} ilustra la salida de la demostración sobre dos ejemplos reales de MAIRA-2:
un informe con una alucinación —una fabricación de infiltrado— que el sistema sitúa en zona roja
(baja confianza), y un informe normal correcto que sitúa en zona verde.

\begin{figure}[t]\centering
\includegraphics[width=0.47\linewidth]{fig_demo_halluc.png}\hfill
\includegraphics[width=0.47\linewidth]{fig_demo_correct.png}
\caption{Salida de la demostración (MAIRA-2) con el anclaje por cajas. Izquierda: informe con una
alucinación de infiltrado, con baja confianza (zona roja). Derecha: informe normal correcto, con alta
confianza (zona verde).}
\label{fig:demo}
\end{figure}

\subsection{El radiólogo como decisor y la mitigación del sesgo de automatización}
El sistema se concibe explícitamente como \textbf{asistivo, no autónomo}: no emite un diagnóstico
final ni bloquea informes, sino que aporta evidencia (anclaje) y una estimación de fiabilidad para que
el \textbf{radiólogo decida}. Este diseño mitiga el \textbf{sesgo de automatización} —la tendencia a
aceptar acríticamente la salida de un sistema automático— de dos formas. Primera, el anclaje se
presenta como puntero de \emph{dónde} mirar, no como prueba de acierto, evitando que una explicación
plausible induzca confianza injustificada. Segunda, la leyenda del semáforo incorpora una
\textbf{advertencia honesta}: el color mide la duda del modelo, de modo que el verde indica baja
probabilidad de \emph{alucinación} pero \emph{no} garantiza la ausencia de \emph{omisiones}, que el
modelo comete con confianza. Comunicar esta limitación en la propia interfaz es una salvaguarda de
seguridad, no un detalle menor. Finalmente, la calibración es \textbf{específica de cada modelo}
—cada generador dispone de su propia calibración del semáforo—, lo que reconoce que la relación entre
confianza y acierto no es universal y debe establecerse por sistema antes de su despliegue.

\section{Consideraciones éticas, legales y de sostenibilidad}
El despliegue de un sistema como el propuesto en un entorno sanitario real está sujeto a un marco
ético y normativo exigente que conviene explicitar. Desde el punto de vista \textbf{legal y
regulatorio}, un sistema que asiste al diagnóstico se considera \emph{producto sanitario} y queda
sujeto a la regulación europea correspondiente (Reglamento de Productos Sanitarios) y, de forma
creciente, al Reglamento de Inteligencia Artificial de la Unión Europea, que clasifica los sistemas de
apoyo a la decisión clínica como de \textbf{alto riesgo} y les impone requisitos de transparencia,
supervisión humana, gestión de riesgos y trazabilidad. El diseño asistivo del presente trabajo —con el
radiólogo como decisor y una salida de confianza explicable y calibrada— se alinea con estos
principios, especialmente con el de \textbf{supervisión humana significativa}. En cuanto a la
\textbf{protección de datos}, el trabajo emplea exclusivamente conjuntos de acceso libre y
anonimizados, evitando el tratamiento de datos personales; un despliegue clínico requeriría, además,
las garantías del Reglamento General de Protección de Datos.

Desde el punto de vista \textbf{ético}, dos riesgos merecen atención. El primero es el \textbf{sesgo}:
un modelo entrenado sobre poblaciones concretas puede rendir de forma desigual en otras, lo que este
trabajo aborda parcialmente al evaluar la generalización a un dominio distinto y al recomendar una
calibración específica por sistema y población antes del despliegue. El segundo es el ya mencionado
\textbf{sesgo de automatización}, cuya mitigación se integra en el propio diseño de la interfaz. En
cuanto a la \textbf{sostenibilidad}, el enfoque prioriza la \textbf{eficiencia}: todos los modelos se
ejecutan cuantizados en una única tarjeta gráfica de consumo, y la señal de fiabilidad propuesta se
obtiene \emph{sin coste adicional de inferencia} —reutiliza las probabilidades que el modelo ya
produce—, evitando el gasto computacional de métodos que exigen múltiples generaciones o modelos
auxiliares. Esta frugalidad favorece la accesibilidad del sistema en entornos con recursos limitados.

\section{Impacto, agentes implicados y viabilidad}
Los \textbf{agentes} implicados en un despliegue del sistema son el \textbf{radiólogo} —usuario y
decisor—, el \textbf{servicio de radiología} —que gestiona la carga de trabajo—, el \textbf{paciente}
—beneficiario último de una lectura más segura— y el \textbf{equipo técnico} responsable de la
calibración y el mantenimiento. El \textbf{impacto} potencial es triple. En el plano
\textbf{asistencial}, un priorizador de revisión bien calibrado permite concentrar la atención del
radiólogo en los informes de mayor riesgo, mitigando la presión asistencial sin comprometer la
seguridad. En el plano de la \textbf{seguridad del paciente}, la comunicación explícita y honesta de
la fiabilidad —incluida la advertencia sobre las omisiones— reduce el riesgo de una confianza
injustificada en la salida automática. Y en el plano \textbf{científico-técnico}, el trabajo aporta
una metodología de auditoría de la fiabilidad reutilizable para otros modelos y dominios.

En cuanto a la \textbf{viabilidad}, el enfoque es deliberadamente \textbf{frugal}: todos los modelos
se ejecutan cuantizados en una única tarjeta gráfica de consumo, sin infraestructura especializada, y
la señal de confianza no incurre en coste de inferencia adicional. El principal recurso consumido es
el \textbf{tiempo de cómputo} de la extracción de señales y el etiquetado, mitigado mediante procesos
reanudables. Esta frugalidad hace el sistema económicamente viable y accesible en entornos con
recursos limitados, y compatible con su integración como módulo asistivo en flujos de trabajo
existentes.

\section{Discusión y análisis de los resultados}
Los resultados dibujan una imagen coherente. La explicación visual de estos modelos es \emph{fiel y
está anclada}, lo que valida su uso como puntero de ``dónde mirar''; pero la fiabilidad \textbf{no se
lee en la forma de la explicación} —ni en el mapa de calor, ni en el \emph{lift} objetivo, ni en las
cajas entrenadas del propio modelo— sino en la \textbf{incertidumbre de la decisión}, y de forma
robusta solo para las alucinaciones. Esta señal, una vez depurada del token de apertura, resulta
\emph{simple, agnóstica al modelo y disponible en inferencia}, y —a diferencia de lo que sugería el
único generador inicial— \emph{se transfiere fuera de distribución} cuando el generador es de calidad
suficiente. Que un modelo mejor mantenga la señal fuera de distribución (0,81) mientras el débil la
pierde (0,61) indica que el problema no es el método sino el generador que, sin datos del dominio,
alucina con confianza. La exploración de señales adicionales (entropía, verosimilitud,
auto-consistencia, dependencia de imagen) mostró que ninguna mejora de forma consistente en los tres
datasets: son específicas de modelo o de dominio, por lo que la fórmula final privilegia la
\emph{simplicidad robusta} de una única señal frente a un combinado sobreajustado.

Las omisiones quedan como problema de \textbf{sensibilidad} del generador, detectable solo en parte
desde sus propias señales y por otra vía (mejor generador o anclaje explícito). Como aportación
transversal, de carácter metodológico, conviene subrayar que repetidamente señales de apariencia
fuerte (el máximo de atención, la entropía cruda confundida por la longitud, la relación
\emph{lift}--acierto, las propias cajas del modelo) resultaron artefactos que solo un control
cuidadoso de confusores y la validación contra referencias objetivas revelaron; en este dominio,
\emph{validar contra verdad objetiva es imprescindible}.

En relación con la literatura, estos resultados matizan y extienden el conocimiento previo. Confirman,
por un lado, la intuición de que la incertidumbre del modelo es un mejor predictor de la alucinación
que sus explicaciones visuales, en línea con los métodos de auto-consistencia y de estimación de
incertidumbre desarrollados para modelos de lenguaje generales; pero aportan, por otro, tres elementos
propios del dominio radiológico que la literatura no había establecido de forma conjunta: la
\textbf{disociación} cuantitativa entre alucinaciones detectables y omisiones que no lo son, la
\textbf{robustez cruzada} de la señal entre modelos y dominios —con la identificación del generador
débil como causa de la degradación fuera de distribución— y la constatación, con las cajas de anclaje
del propio modelo, de que ni siquiera una explicación entrenada y plausible predice el acierto. El
valor práctico de estos hallazgos es directo: indican que un sistema de ayuda al radiólogo debe
construir su señal de confianza sobre la incertidumbre de la decisión, comunicar honestamente su
alcance (que no cubre las omisiones) y calibrarse por modelo antes de su despliegue.

Finalmente, procede señalar las limitaciones:
MAIRA-2 fue entrenado con PadChest, por lo que su evaluación in-distribution puede incluir solape
train--test —mitigado usando Open-I, ajeno a su entrenamiento, como prueba fuera de distribución—; el
número de alucinaciones por celda es reducido, lo que ensancha los intervalos de confianza; y el
beneficio clínico real de presentar informe, anclaje y confianza al radiólogo requeriría un estudio
prospectivo de lectura, que se propone como trabajo futuro.

\chapter{Conclusiones}
% BORRADOR autónomo (2026-08-08). Editar a tu voz. UAX: empezar por objetivos específicos y
% terminar por el general, mostrando su cumplimiento. Recomendado: 5-10 páginas.

Este capítulo recupera los objetivos planteados y valora su cumplimiento, recorre las limitaciones
y la prospectiva, y cierra con una reflexión personal.

Respecto a los \textbf{objetivos específicos}, todos se han cumplido. Se ha \textbf{revisado el
estado del arte} (Cap.~3), situando la aportación en el hueco de la explicabilidad validada de los
MLLM de radiografía de tórax. Se ha \textbf{diseñado e implementado} un generador multimodal propio
con destilación de conocimiento y un marco reutilizable de explicabilidad. Se ha \textbf{cuantificado
de forma objetiva el anclaje clínico} de la atención contra anotaciones reales de localización
(lift global 3,03). Y se ha \textbf{determinado la capacidad de la explicabilidad y la incertidumbre
para detectar errores}, con un resultado matizado y sólido: la morfología de la explicación no
predice la fiabilidad —lo que se confirma por triple vía, incluidas las cajas de anclaje del propio
modelo—, pero la \textbf{incertidumbre de la decisión detecta las alucinaciones} (AUC 0,90 tras
depurar el token de apertura), no las omisiones, que solo se captan de forma parcial.

En cuanto al \textbf{objetivo general} —evaluar la fidelidad y el anclaje de las explicaciones y en
qué medida permiten estimar la fiabilidad—, se ha alcanzado y ampliado: las explicaciones son fieles
y ancladas, pero la fiabilidad reside en la incertidumbre de la decisión, no en el mapa de calor. Es
más, dicha señal ha demostrado ser \textbf{agnóstica al modelo y robusta al dominio}: validada en un
diseño \(2\times2\) (dos generadores \(\times\) dos dominios), se transfiere a un segundo modelo del
estado del arte (MAIRA-2) y se mantiene \emph{fuera de distribución} para el generador de calidad
suficiente (AUC 0,80--0,81), mientras que la degradación observada se atribuye al generador débil, no
al método. La contribución es, por tanto, triple: la \textbf{auditoría objetiva del anclaje} de un
MLLM de CXR; la \textbf{disociación} entre errores detectables (fabricaciones) e indetectables por
confianza (omisiones); y una \textbf{señal de confianza simple, agnóstica, calibrada y validada de
forma cruzada}, junto con su traducción clínica en una demostración interactiva. Como aportación
metodológica menor pero de valor práctico, se documenta que \emph{excluir el token de apertura del
informe} mejora de forma sistemática la señal de incertidumbre.

El significado de conjunto de estas conclusiones puede resumirse en una idea: en los modelos
multimodales de radiografía de tórax, la fiabilidad de un informe no se lee en \emph{lo que el modelo
muestra} —su explicación visual— sino en \emph{lo seguro que está} al decidir, y esa seguridad es
cuantificable, transferible entre modelos y calibrable para su uso clínico, siempre que se comunique
con honestidad su alcance. Este desplazamiento del foco —de la explicación a la incertidumbre— es la
contribución conceptual central del trabajo, y tiene implicaciones directas para el diseño de sistemas
de ayuda a la decisión: la evidencia visual debe presentarse como \emph{apoyo a la localización}, no
como garantía de acierto, y la confianza debe construirse sobre la señal interna del modelo, no sobre
la forma de su explicación.

\clearpage
\section{Respuesta a las preguntas de investigación}

\begin{table}[h]\centering
\caption{Resumen de las cuatro preguntas de investigación (Capítulo~1) e hipótesis asociadas
(H1--H4, Sección~4.2) frente a la evidencia obtenida.}
\label{tab:resumen_pi}
\begin{tabular}{lp{2.6cm}p{5.5cm}}
\toprule
Pregunta & Resultado & Evidencia clave \\ \midrule
PI1 — ¿anclaje clínico? & Confirmada & Lift global 3,03; \textit{pointing-game} 0,38 \\
PI2 — ¿morfología predice fiabilidad? & Refutada & AUC 0,54 (n.s.); confirmado por triple vía \\
PI3 — ¿incertidumbre detecta errores? & Confirmada (asimétrica) & AUC 0,80--0,90 alucinaciones; parcial en omisiones \\
PI4 — ¿agnóstica y robusta? & Confirmada con matiz & Transfiere a MAIRA-2; falla solo en generador débil OOD \\
\bottomrule
\end{tabular}
\end{table}

Recapitulando las cuatro preguntas planteadas en el Capítulo~1: en cuanto a \textbf{PI1}, la atención
está clínicamente anclada (lift 3,03; tres veces el azar en el \textit{pointing-game}), lo que valida
su uso como puntero de \emph{dónde} mira el modelo. Respecto a \textbf{PI2}, la morfología de la
explicación \emph{no} predice la fiabilidad: ni el mapa de calor de atención, ni el \textit{lift}
objetivo, ni las cajas de anclaje entrenadas del propio modelo —que localizan las alucinaciones con
la misma plausibilidad que los hallazgos reales— separan el acierto del error. En cuanto a
\textbf{PI3}, la incertidumbre de la decisión sí detecta los errores, pero de forma \emph{asimétrica}:
robusta para las alucinaciones (AUC 0,80--0,90) y solo parcial para las omisiones, que tienden a
situarse cerca del umbral. Y sobre \textbf{PI4}, la señal es agnóstica al modelo y robusta al dominio:
se transfiere a un segundo generador y se mantiene fuera de distribución para un generador de calidad
suficiente, degradándose solo con un generador débil que, sin datos del dominio, alucina con
confianza. En conjunto, las cuatro preguntas quedan respondidas de forma coherente y con matices
respaldados por la evidencia.

\section{Limitaciones}
Algunas limitaciones de las versiones preliminares se han \emph{resuelto} en el propio trabajo: la
dependencia de un único juez se ha mitigado con \textbf{tres fuentes de verdad} concordantes; la de un
único modelo, validando en \textbf{dos generadores}; y la de un único dominio, con la \textbf{réplica
en PadChest-GR}. Persisten, no obstante, limitaciones reales. El número de alucinaciones por celda es
\textbf{reducido} (entre 12 y 110), lo que ensancha los intervalos de confianza. La evaluación
\emph{in-distribution} de MAIRA-2 sobre PadChest puede incluir \textbf{solape train--test}, al haber
sido entrenado con ese corpus; se mitiga usando Open-I —ajeno a su entrenamiento— como prueba fuera de
distribución, pero conviene señalarlo. La detección de \textbf{omisiones} es solo parcial: son errores
de sensibilidad del generador, no siempre acompañados de incertidumbre. Metodológicamente, varias
señales de forma resultaron engañosas (el máximo de atención, la entropía confundida con la longitud,
la relación \emph{lift}--acierto), lo que obligó a un control cuidadoso de confusores. Por último, el
\textbf{beneficio clínico no se ha demostrado empíricamente}: requeriría un estudio prospectivo de
lectura.

Varias de estas limitaciones tienen origen en las condiciones materiales del proyecto. El
\textbf{acceso a datos} clínicos de mayor tamaño estuvo restringido: MIMIC-CXR, el corpus narrativo
de referencia en el área, exige credenciales de PhysioNet no disponibles para este trabajo
(Sección~4.1), lo que obligó a construir el diseño experimental sobre conjuntos abiertos de menor
tamaño (Open-I, NIH ChestX-ray8, PadChest-GR) y es, en buena medida, la causa directa del reducido
número de alucinaciones por celda señalado más arriba. El \textbf{cómputo} también estuvo acotado a
una única GPU de consumo, lo que impuso ejecutar los modelos del estado del arte cuantizados a 8
bits —sin comparación frente a precisión completa— y limitó la escala y el número de repeticiones
de las ablaciones del modelo propio (Sección~4.4). Ninguna de estas restricciones invalida los
resultados, que se sostienen sobre tres fuentes de verdad y dos réplicas independientes, pero sí
marca el techo de precisión alcanzable con los recursos disponibles y es la razón principal de la
primera línea de prospectiva (consolidación estadística).

\section{Prospectiva}
El trabajo abre varias líneas de continuación, que se ordenan de la más inmediata a la más ambiciosa.

\textbf{Consolidación estadística.} La primera línea es \textbf{ampliar la muestra} de imágenes por
celda del diseño, dado que el número de alucinaciones —el factor que limita la precisión de la
estimación— es reducido; ampliarlo estrecharía los intervalos de confianza. En la misma dirección,
añadir más generadores (por ejemplo, modelos con razonamiento por agentes o de otra familia) y más
dominios reforzaría la generalidad de la conclusión sobre la robustez de la señal.

\textbf{Cierre de la duda del solape.} La evaluación \emph{in-distribution} de MAIRA-2 sobre PadChest
podría incluir solape con su entrenamiento; evaluar el modelo fuera de distribución en dominios
verdaderamente ajenos y sin intersección con su corpus cerraría definitivamente esta cuestión.

\textbf{Detección de omisiones.} Las omisiones son el error que la incertidumbre no capta de forma
fiable, por tratarse de fallos confiados de sensibilidad. Abordarlas requeriría vías complementarias:
mejorar la sensibilidad del propio generador, emplear un anclaje explícito que ``obligue'' a mirar
regiones sospechosas, o recurrir a la \textbf{auto-consistencia semántica} entre múltiples muestras,
que podría revelar la inestabilidad de un informe incompleto.

\textbf{De la detección a la prevención.} Más allá de detectar la alucinación tras generarla, la señal
de incertidumbre podría integrarse en el propio proceso de \emph{decodificación} para \textbf{prevenir}
la fabricación —por ejemplo, penalizando las continuaciones de baja confianza o solicitando evidencia
adicional—, convirtiendo el sistema de un detector \emph{a posteriori} en una salvaguarda \emph{a
priori}.

\textbf{Validación clínica.} Finalmente, la línea de mayor calado es un \textbf{estudio de lectura}
prospectivo con radiólogos que mida el beneficio real de presentar informe, anclaje y confianza,
cuantificando su efecto sobre la precisión diagnóstica, el tiempo de lectura y el riesgo de sesgo de
automatización. Solo un estudio de este tipo puede demostrar el valor clínico que este trabajo
fundamenta pero no llega a probar, y sería el paso natural hacia una eventual integración del sistema
como módulo asistivo en el flujo de trabajo radiológico.

\clearpage
\section{Mejoras futuras del sistema}
Las líneas anteriores extienden la investigación; las que siguen extienden el \textbf{sistema} en un
sentido de ingeniería, orientado a acercarlo a un entorno de uso real.

\textbf{Ampliación del catálogo de hallazgos.} El etiquetado objetivo se restringe a los hallazgos
\emph{agudos} más frecuentes y a las ocho patologías con caja de NIH ChestX-ray8. Extender el marco a
un catálogo más amplio —hallazgos crónicos, óseos o de dispositivos— exigiría nuevas fuentes de verdad,
pero ampliaría la cobertura clínica real del sistema.

\textbf{Integración con el flujo de trabajo radiológico.} La demostración actual es un prototipo
\textit{standalone}; su valor práctico aumentaría integrándola como módulo dentro del visor de
imágenes o del sistema de información radiológica (RIS/PACS) del centro, de modo que el informe, el
anclaje y la confianza aparezcan directamente en la estación de lectura, sin cambiar la herramienta que
el radiólogo ya usa.

\textbf{Optimización de la latencia de inferencia.} El sistema se ejecuta hoy cuantizado a 8 bits en
una única GPU de consumo (Sección~4.7); técnicas adicionales de optimización —cuantización a 4 bits,
motores de inferencia especializados o destilación del propio generador— reducirían la latencia y el
coste de despliegue, acercando el sistema a un uso en tiempo real durante la lectura.

\textbf{Ampliación a otras modalidades.} El marco de explicabilidad e incertidumbre —agnóstico al
generador por construcción— no depende de particularidades de la radiografía de tórax; extenderlo a
otras modalidades con MLLM propios (tomografía computarizada, resonancia magnética, mamografía)
permitiría comprobar si la conclusión central del trabajo —que la fiabilidad reside en la incertidumbre
y no en la explicación visual— se sostiene fuera del dominio aquí estudiado.

\textbf{Modelos de mayor escala y capacidad de imagen.} Este trabajo se ha ejecutado sobre generadores
de 7 mil millones de parámetros y con la resolución de entrada que cada uno define por diseño,
condicionado por el hardware de consumo disponible (Sección~6.2). Investigar con modelos de mayor
potencia y con mayor capacidad de procesamiento de imagen —más tokens visuales, mayor resolución de
entrada— permitiría estudiar si la morfología del mapa de calor, que en este trabajo no logró predecir
la fiabilidad, se comporta de forma distinta en arquitecturas más grandes y con una representación
visual más rica, cerrando una pregunta que la limitación de recursos deja abierta.

\textbf{Consolidación de la portabilidad.} El trabajo ya distribuye la memoria y las demostraciones
mediante imágenes \textbf{Docker} reproducibles (Anexo de organización del repositorio); una línea
natural es exponer el sistema de análisis como una \textbf{API} desplegable en la nube, desacoplada de
la interfaz Gradio actual, para que otros equipos puedan integrarlo en sus propias herramientas sin
depender de esta implementación concreta.

\section{Consideraciones finales}
Considero que el resultado final de esta investigación fue satisfactorio, no solo por lo logrado al
poder ``acotar'' con una métrica un umbral de alucinaciones —con el beneficio añadido de que muchas
omisiones se sitúan precisamente en ese umbral—, sino por todo el aprendizaje que ha conllevado.

Ese aprendizaje me ha enseñado, incluso creyendo que ya estaba familiarizado con el método científico,
lo que de verdad significa el rigor en el control de confusores tras varios ``falsos positivos''.
Cuántas veces habré creído tener algo interesante entre manos, para que la siguiente batería de
pruebas mostrara que la hipótesis no se sostenía. Pero esa negación de la idea inicial del TFM no fue
el final del proyecto, sino el paso de una intuición —los mapas de calor— a su refutación con
evidencia. He aprendido a trabajar con MLLM de gran escala, con cuantización y con extracción de
señales a escala; a comunicar un resultado negativo con la misma solidez que uno positivo; y,
finalmente, a llegar a algo que sí era capaz de aportar.

Si hoy empezara de nuevo —no con el conocimiento adquirido, sino con las aptitudes desarrolladas
durante el Máster y el propio TFM— seguramente no habría cambiado el resultado, pues se habrían hecho
las mismas pruebas, pero las habría hecho bien a la primera, con más rigor científico y sin dejarme
llevar por pequeñas victorias sesgadas que acabaron en sobretrabajo antes de comprobar si realmente
eran buenos indicios. Y, sobre todo, habría empezado el procesamiento de datos mucho antes, incluso
mientras todavía planteaba las ideas. Me habría encantado poder validar las hipótesis con más datos,
pero el tiempo disponible y el hardware del equipo con el que se realizó el trabajo lo complicaron.

No puedo terminar este trabajo sin expresar mi deseo de que la asistencia por IA en radiología
evolucione hacia sistemas honestos sobre su propia incertidumbre, y no solo hacia modelos cada vez más
grandes —y, en muchos casos, aún más caja negra—, y de que se apliquen a los MLLM todos los avances de
la IA explicable clásica, como los vistos en la asignatura de XAI del Máster. Espero que esta pequeña
contribución sea un paso, por modesto que sea, en esa dirección.


% --- Bibliografía ---
\chapter*{Bibliografía}\addcontentsline{toc}{chapter}{Bibliografía}
\printbibliography[heading=none]

% --- Anexos (formato libre, no puntúan, máx. ~20% del contenido) ---
\appendix
\chapter{Anexos: detalles de reproducibilidad}

Este anexo recoge los detalles técnicos que permiten reproducir los experimentos. No forma parte de
la evaluación pero documenta las decisiones de implementación.

\section{Entornos de ejecución}
Los experimentos se ejecutan en Python sobre dos entornos aislados, dado que los dos generadores
exigen versiones incompatibles de la biblioteca de \textit{transformers}. El entorno de LLaVA-Rad
emplea \texttt{transformers 4.31} y \texttt{torch 2.0}; el de MAIRA-2, \texttt{transformers 4.46} y
\texttt{torch 2.1}. Ambos generadores se cargan con cuantización a 8 bits mediante
\texttt{bitsandbytes}, lo que permite ejecutar modelos de 7\,mil millones de parámetros en una única
tarjeta gráfica de 12\,GB de memoria.

\section{Extracción de señales}
Para cada imagen se realiza una pasada que produce: (i) el informe voraz junto con las probabilidades
y entropías por token; (ii) $K=4$ muestras estocásticas para la auto-consistencia; y (iii) una
generación con la imagen en blanco para la dependencia de imagen. En el subconjunto con mapa de calor
se añade una pasada con salida de atención. La extracción es resumible: los resultados se
escriben incrementalmente en formato JSONL, de modo que un proceso interrumpido se reanuda saltando
los identificadores ya procesados. El rendimiento es de aproximadamente 27--42\,segundos por imagen
según el modelo.

\section{Hiperparámetros de generación}
El informe principal se genera de forma voraz (\textit{greedy}, sin muestreo) con un máximo
de 60--120 tokens nuevos. Las $K$ muestras de auto-consistencia se generan con muestreo estocástico
(temperatura $0{,}8$, \textit{top-}$p$ $0{,}95$). La generación con imagen en blanco emplea una imagen
de valor cero del mismo tamaño. Para MAIRA-2, el anclaje se obtiene activando la generación con cajas
y analizando la secuencia resultante en frases y coordenadas normalizadas.

\section{Acceso a los datos y etiquetado}
Se emplean tres conjuntos de acceso libre: Open-I / Indiana University (informes y términos MeSH), NIH
ChestX-ray8 (cajas de ocho patologías) y un subconjunto de PadChest-GR (hallazgos presentes/ausentes
y cajas). El etiquetado por juez se ejecuta de forma reproducible y con escritura incremental, con una
rúbrica fija que restringe la evaluación a los hallazgos agudos significativos.

\section{Métodos estadísticos}
La capacidad discriminativa se mide con el área bajo la curva ROC, calculada mediante el estadístico
de Mann--Whitney sobre rangos para garantizar reproducibilidad. Los intervalos de confianza se
estiman por \textit{bootstrap} estratificado (2000 remuestreos) y la significancia frente al azar, con
tests de permutación. La calibración se evalúa por bandas de riesgo definidas sobre cuartiles de la
distribución de los informes correctos. Todo el control de confusores (longitud, posición del token,
tipo de mapa de calor) se realiza mediante residualización o comparaciones emparejadas.

\section{Hardware}
Todos los experimentos se ejecutan localmente en una tarjeta gráfica NVIDIA de gama de consumo con
12\,GB de memoria, lo que demuestra la viabilidad del enfoque sin infraestructura especializada.

\chapter{Anexos: manual de usuario de la demostración}

Este anexo describe el uso de la demostración interactiva desarrollada, que materializa el sistema de
informe, anclaje y confianza. Se proporcionan dos aplicaciones equivalentes, una por generador, dado
que ambos modelos no caben simultáneamente en la memoria de una única tarjeta gráfica de 12\,GB.

\section{Puesta en marcha}
Cada aplicación se lanza de forma independiente desde su entorno correspondiente. La aplicación de
LLaVA-Rad se sirve en un puerto local y muestra el anclaje mediante mapa de calor; la de
MAIRA-2, en otro puerto, muestra el anclaje mediante cajas generadas por el propio modelo.
Al arrancar, la aplicación carga el modelo cuantizado (unos treinta segundos) y queda accesible desde
un navegador web en la dirección local indicada. Alternativamente, cada aplicación puede lanzarse
mediante la imagen Docker correspondiente (Anexo de organización del repositorio), lo que
evita instalar manualmente el entorno Python y sus versiones fijadas; en ese caso, los pesos del
modelo se descargan automáticamente en el primer arranque.

\section{Interfaz y flujo de trabajo}
La interfaz se organiza en tres paneles, alineados con la propuesta del trabajo:
\begin{itemize}
  \item \textbf{Panel izquierdo: informe.} Muestra la radiografía cargada y, debajo, el informe
  generado por el modelo. El usuario puede subir su propia imagen o seleccionar uno de los ejemplos
  incluidos, que además presentan el \emph{veredicto real del radiólogo} para comparar.
  \item \textbf{Panel central: anclaje.} Muestra \emph{dónde} mira el modelo: el mapa de calor de
  atención (LLaVA-Rad) o las cajas de localización (MAIRA-2). Sirve como evidencia visual de la
  afirmación, no como prueba de acierto.
  \item \textbf{Panel derecho: confianza.} Muestra el semáforo calibrado: una banda de
  verde a rojo con un marcador en la posición de la puntuación del informe, la banda de riesgo
  correspondiente y una tabla con la probabilidad de alucinación de cada banda. Incluye la advertencia
  de que el verde indica baja probabilidad de alucinación pero no garantiza la ausencia de omisiones.
\end{itemize}
El flujo de uso previsto es: cargar la imagen, pulsar \emph{Analizar} y leer el informe consultando el
anclaje como apoyo y el semáforo para priorizar la revisión de los casos en zona roja. Las
Figuras~\ref{fig:app_llavarad} y~\ref{fig:app_maira} muestran las dos aplicaciones en ejecución sobre
un caso real de cada conjunto.

\begin{figure}[p]\centering
\includegraphics[width=0.95\linewidth]{fig_app_llavarad.png}
\caption{Aplicación de LLaVA-Rad en ejecución sobre un caso real de Open-I. El modelo alucina una
masa pulmonar inexistente (compárese con el veredicto real del radiólogo, panel superior); el
semáforo lo detecta correctamente y lo sitúa en zona roja (\emph{min\_prob} 0,11, P(alucinación)
$\approx$ 30\,\%).}
\label{fig:app_llavarad}
\end{figure}

\begin{figure}[p]\centering
\includegraphics[width=0.95\linewidth]{fig_app_maira.png}
\caption{Aplicación de MAIRA-2 en ejecución sobre un caso real de PadChest-GR. El modelo también
alucina en este caso; el semáforo lo sitúa igualmente en zona roja (\emph{min\_prob} 0,12,
P(alucinación) $\approx$ 14\,\%). El panel de anclaje muestra la imagen con las cajas propias del
modelo cuando las genera.}
\label{fig:app_maira}
\end{figure}
\FloatBarrier

\section{Instalación y ejecución con Docker (paso a paso)}
Esta sección explica cómo poner en marcha ambas demostraciones en un ordenador que no tiene instalado
ni Python ni las bibliotecas del proyecto (por ejemplo, el equipo de un profesor evaluador), usando
únicamente Docker. No se requiere ningún conocimiento de programación; basta con ejecutar los
comandos indicados en una terminal.

\paragraph{Requisitos.} Un equipo con Docker Desktop instalado (gratuito, disponible para
Windows, macOS y Linux desde \url{https://www.docker.com/products/docker-desktop/}) y, para las
demostraciones (no para compilar la memoria), una tarjeta gráfica NVIDIA con al menos
12\,GB de memoria y el \textit{NVIDIA Container Toolkit} instalado, que permite a Docker acceder a la
GPU. Sin GPU es posible, aun así, compilar la memoria en PDF con la primera imagen descrita a
continuación.

\paragraph{Paso 1: obtener el código.} Descomprimir el paquete del proyecto (o clonar el repositorio)
en cualquier carpeta local, y abrir una terminal en la carpeta raíz del proyecto (la que contiene la
carpeta \texttt{docker/}).

\paragraph{Paso 2: construir las imágenes.} Cada imagen se construye una única vez; la construcción
tarda varios minutos porque descarga las bibliotecas necesarias, pero no requiere ninguna interacción:
\begin{verbatim}
docker build -t tfm-llavarad -f docker/Dockerfile.llavarad .
docker build -t tfm-maira2   -f docker/Dockerfile.maira2   .
\end{verbatim}

\paragraph{Paso 3: lanzar una demostración.} Las dos aplicaciones no caben a la vez en una GPU de
12\,GB, por lo que se lanzan una cada vez:
\begin{verbatim}
docker run --rm --gpus all -p 7860:7860 \
  -v hf_cache_llavarad:/hf_cache tfm-llavarad
docker run --rm --gpus all -p 7862:7862 \
  -v hf_cache_maira2:/hf_cache tfm-maira2
\end{verbatim}
La primera vez que se ejecuta cada comando, el contenedor descarga automáticamente los pesos del
modelo desde \textit{Hugging Face Hub} (varios minutos, según la conexión) y los conserva en el
volumen indicado, de modo que los arranques posteriores son mucho más rápidos. El acceso al modelo
\texttt{microsoft/maira-2} exige haber aceptado su licencia en Hugging Face con la cuenta cuyo token
se utilice; el contenedor solicitará este dato la primera vez si no está ya configurado mediante la
variable de entorno \texttt{HF\_TOKEN}.

\paragraph{Paso 4: abrir la demostración.} Cuando la terminal muestra que el servidor está listo,
basta con abrir un navegador web en \url{http://127.0.0.1:7860} (LLaVA-Rad) o
\url{http://127.0.0.1:7862} (MAIRA-2) y usar la interfaz descrita en la sección anterior. Para
detener la aplicación es suficiente con cerrar la terminal o pulsar \texttt{Ctrl+C}.

\paragraph{Compilar la memoria sin instalar LaTeX.} De forma análoga, la memoria en PDF puede
regenerarse sin instalar ninguna distribución LaTeX localmente:
\begin{verbatim}
docker build -t tfm-latex -f docker/Dockerfile.latex .
docker run --rm \
  -v "RUTA_ABSOLUTA_AL_PROYECTO/memoria:/memoria" tfm-latex
\end{verbatim}
sustituyendo \texttt{RUTA\_ABSOLUTA\_AL\_PROYECTO} por la ruta local del proyecto; el PDF resultante
se escribe directamente en la carpeta \texttt{memoria/}. Los detalles adicionales (versiones fijadas,
\texttt{docker compose} como alternativa) se documentan en el Anexo de organización del repositorio.

\section{Interpretación de la salida}
Un informe en zona verde indica que el modelo ha generado su veredicto con alta confianza y
baja probabilidad de alucinación; el radiólogo puede revisarlo con menor prioridad, sin olvidar la
posibilidad de una omisión. Un informe en zona roja señala alta incertidumbre y una
probabilidad de alucinación elevada; debe revisarse de forma prioritaria. Las bandas intermedias
gradúan esa prioridad. En todo caso, la decisión final corresponde siempre al facultativo.

\chapter{Anexos: organización del repositorio}

El código del proyecto se organiza de forma modular para separar responsabilidades y facilitar la
reproducibilidad, siguiendo la arquitectura por capas descrita en el Capítulo~4.

\section{Estructura}
El módulo de datos agrupa la ingesta y el preprocesado de los conjuntos empleados. El módulo
de modelos contiene la arquitectura del generador propio (codificador visual, decodificador
con atención cruzada y su ensamblaje). El módulo de explicabilidad implementa la extracción
del mapa de calor, las métricas de anclaje y la evaluación de fidelidad por oclusión. El módulo de
entrenamiento orquesta las dos etapas (alineamiento contrastivo y ajuste generativo). Los
scripts de extracción producen, para cada modelo y conjunto, los ficheros de señales en
formato JSONL reanudable, y los scripts de análisis calculan las métricas de detección y calibración.
Finalmente, la demostración contiene las dos aplicaciones interactivas.

\section{Reproducibilidad}
Cada experimento se acompaña de sus scripts y de los ficheros intermedios (señales, etiquetas,
figuras), de modo que los resultados de la memoria pueden regenerarse. Los procesos largos se ejecutan
de forma reanudable, escribiendo su salida de forma incremental, lo que permite reanudarlos tras una
interrupción sin repetir el trabajo ya realizado.

\section{Contenerización y portabilidad}
Para evitar que la reproducibilidad dependa de instalar manualmente una distribución LaTeX o los dos
entornos Python de inferencia (cada uno con una combinación de versiones específica y ya
depurada de conflictos entre \textit{bibliotecas}), el repositorio incluye tres imágenes
Docker independientes. La primera compila la memoria (\texttt{pdflatex} + \texttt{biber} +
\texttt{biblatex-apa} + TikZ) sin requerir ninguna instalación local de LaTeX; basta con montar el
directorio \texttt{memoria/} como volumen. Las otras dos reproducen, respectivamente, el entorno de
LLaVA-Rad y el de MAIRA-2 empleados en las demostraciones (Anexo de manual de usuario), fijando las
mismas versiones validadas en el Anexo de reproducibilidad (incluida la combinación específica
de \texttt{gradio~3.50.2} con \texttt{pydantic~1.x} necesaria para evitar la incompatibilidad con
\texttt{pydantic~2} presente en versiones más recientes de Gradio). Los pesos de los modelos no
se empaquetan en las imágenes, que quedarían así de decenas de gigabytes: se descargan de
\textit{Hugging Face Hub} en el primer arranque de cada contenedor y se conservan en un volumen
persistente para no repetir la descarga.

Como el objetivo tras concluir la fase de entrenamiento es poder continuar la redacción de la memoria
en cualquier dispositivo sin necesidad de GPU, el repositorio se distribuye también en un
paquete portable reducido (código, memoria, demostraciones y los resultados intermedios del
análisis) que excluye deliberadamente los conjuntos de datos crudos, los \textit{checkpoints} de
entrenamiento y la caché de modelos (del orden de 90\,GB en conjunto y prescindibles para este fin),
documentando en un fichero de nivel superior cómo y por qué se han omitido, y cómo recuperarlos si en
el futuro fuese necesario reanudar el entrenamiento.

\end{document}
